% Informe del Simulador Inmersivo de Boxeo en RV (1P vs. IA)
% Entregable 5 (semana 14): versión final, 100 % de avance acumulado.
% Compilar con: xelatex (dos pasadas) o latexmk -xelatex
% Lo que depende del visor y de las pruebas con usuarios está marcado con \pendiente{...} (en rojo); la portada dice
% cuántos quedan. No se entrega mientras quede alguno.
\documentclass[12pt,letterpaper]{article}

\usepackage{fontspec}
\setmainfont{texgyretermes}[Extension=.otf,UprightFont=*-regular,BoldFont=*-bold,ItalicFont=*-italic,BoldItalicFont=*-bolditalic]
\setmonofont{texgyrecursor}[Extension=.otf,UprightFont=*-regular,BoldFont=*-bold,ItalicFont=*-italic,BoldItalicFont=*-bolditalic,Scale=0.88]
\usepackage[spanish,es-tabla,es-noshorthands]{babel}
\usepackage[top=2.5cm,bottom=2.5cm,left=2.5cm,right=2.5cm]{geometry}
\usepackage{graphicx}
\usepackage{xcolor}
\usepackage{colortbl}
\usepackage{array}
\usepackage{tabularx}
\usepackage{multirow}
\usepackage{makecell}
\usepackage{longtable}
\usepackage{enumitem}
\usepackage{titlesec}
\usepackage{fancyhdr}
\usepackage{caption}
\usepackage{subcaption}
\usepackage{float}
\usepackage{needspace}
\usepackage{microtype}
\usepackage{tikz}
\usepackage{xurl}
\usepackage[hidelinks,bookmarksnumbered]{hyperref}

% ---- colores (los del cronograma del informe original) ----
\definecolor{vino}{HTML}{8B0A1E}
\definecolor{rojo}{HTML}{C8102E}
\definecolor{dorado}{HTML}{D4A62A}
\definecolor{crema}{HTML}{F2E4B4}
\definecolor{gris}{HTML}{E6E6E6}

\graphicspath{{img/}{../capturas/}}
\setlength{\parindent}{0pt}
\setlength{\parskip}{6pt}
\renewcommand{\arraystretch}{1.18}
\setlist{itemsep=2pt,topsep=3pt}
\setlist[itemize]{label=\textbullet}
\emergencystretch=2em

\captionsetup{font=small,labelfont=bf,skip=4pt}

% ---- encabezados de sección ----
\titleformat{\section}{\Large\bfseries}{\thesection.}{0.5em}{}
\titleformat{\subsection}{\large\bfseries}{\thesubsection}{0.5em}{}
\titleformat{\subsubsection}{\normalsize\bfseries}{\thesubsubsection}{0.5em}{}
\titlespacing*{\section}{0pt}{16pt}{6pt}
\titlespacing*{\subsection}{0pt}{12pt}{4pt}
\titlespacing*{\subsubsection}{0pt}{8pt}{3pt}

\pagestyle{fancy}
\fancyhf{}
\fancyfoot[C]{\thepage}
\fancyhead[L]{\small\itshape Simulador inmersivo de boxeo en RV (1P vs.\ IA)}
\fancyhead[R]{\small\itshape Entregable 5 (semana 14)}
\renewcommand{\headrulewidth}{0.3pt}
\fancypagestyle{plain}{\fancyhf{}\fancyfoot[C]{\thepage}\renewcommand{\headrulewidth}{0pt}}

% ---- atajos ----
% Nombre de código (identificador, archivo, campo): monoespaciado, no se corta.
\newcommand{\cs}[1]{\texttt{\detokenize{#1}}}
% Encabezado de tabla: fondo vino, texto blanco.
\newcommand{\hd}[1]{\textcolor{white}{\textbf{#1}}}
\newcolumntype{L}{>{\raggedright\arraybackslash}X}
\newcolumntype{C}{>{\centering\arraybackslash}X}
\newcommand{\cab}{\rowcolor{vino}}
\newcommand{\ok}{\textbf{Completado}}
% Dato que falta (visor o pruebas con usuarios): visible en rojo, cuenta para la portada y avisa al compilar.
\newcounter{pendientes}
\newcommand{\pendiente}[1]{\stepcounter{pendientes}\textcolor{red}{\textbf{[PENDIENTE: #1]}}\GenericWarning{}{Informe: hay un pendiente sin rellenar}}
\makeatletter
\AtEndDocument{\immediate\write\@auxout{\gdef\string\pendientesTotales{\thependientes}}}
\makeatother

% Tabla estándar: \tabla{especificación}{encabezado}{cuerpo}
\newcommand{\tablita}[3]{%
  \begingroup\small
  \begin{tabularx}{\linewidth}{#1}
    \hline
    \cab #2 \\
    \hline
    #3
    \hline
  \end{tabularx}
  \endgroup}

\begin{document}
\pagenumbering{gobble}

% ================= PORTADA =================
\begin{titlepage}
  \centering
  \vspace*{0.5cm}
  \includegraphics[width=0.55\linewidth]{logo_utec.png}\\[1.6cm]
  {\LARGE\bfseries SIMULADOR INMERSIVO DE BOXEO EN REALIDAD VIRTUAL (1P VS.\ IA)\par}
  \vspace{0.4cm}
  {\small\itshape Diseño de Interacción Espacial, Presencia y Evaluación de Usabilidad en Unity para Meta Quest\par}
  \vspace{1.4cm}
  {\large\textbf{Docente:} Chambilla, Teofilo\par}
  \vspace{0.3cm}
  {\large\textbf{Curso:} Arte y Tecnología\par}
  \vspace{0.3cm}
  {\large\textbf{Entregable 5 (semana 14):} versión final, avance del 100\,\%\par}
  \ifdefined\pendientesTotales\ifnum\pendientesTotales>0
    \vspace{0.3cm}
    {\large\color{red}\bfseries BORRADOR: quedan \pendientesTotales\ datos pendientes (en rojo en el texto)\par}
  \fi\fi
  \vspace{1cm}
  \begin{tabular}{|c|}
    \hline
    \rowcolor{gris}\textbf{Integrantes}\\ \hline
    Silva Reyes, Santiago Miguel\\ \hline
    Roque, Estefano\\ \hline
    Silva Mena, David Manuel\\ \hline
  \end{tabular}\\[1.6cm]
  {\large Lima -- Perú\\ 2026\par}
\end{titlepage}

\pagenumbering{arabic}
\setcounter{page}{2}
\tableofcontents
\clearpage

% Las secciones que cambian respecto a la semana 13 tienen copia propia (_s14), para no alterar ese informe
% ============================================================
\section{Introducción}
Nuestro proyecto consiste en un simulador inmersivo de boxeo en Realidad Virtual enfocado en la modalidad de un jugador contra la inteligencia artificial, desarrollado en el motor Unity para visores Meta Quest. El sistema aborda de manera directa el problema de la monotonía y la ausencia de retroalimentación técnica inmediata al realizar rutinas de ejercicio cardiovascular o entrenamiento de combate en el hogar. La experiencia sitúa al usuario dentro de un cuadrilátero tridimensional a escala real donde debe enfrentar a una serie progresiva de cinco oponentes virtuales organizados en niveles de dificultad ascendente, permitiendo perfeccionar tiempos de reacción, coordinación visomotora y esquives activos en un entorno seguro e inmersivo.

Este documento reúne el diseño del proyecto y el registro de lo implementado en cada entregable. La sección~\ref{sec:e5} detalla el Entregable~5 (semana 14), la versión final: el quinto rival (Iron Mike Tyson), la escena de resultados con el cinturón de campeón, la preparación del rendimiento para Meta Quest~2 y de las pruebas de usabilidad HEART, y la comparación de los bocetos con lo implementado. La sección~\ref{sec:e4} detalla el Entregable~4 (semana 13): cuatro rivales jugables, rondas de tres rounds de 90 segundos, réferi con cuenta regresiva y ciclo de KO, guardias grabadas para los bloqueos, háptica diferenciada completa y registro de métricas por sesión.

% ============================================================
\section{Antecedentes}
El entrenamiento cardiovascular y las disciplinas de combate practicadas en el hogar enfrentan una limitación estructural difícil de resolver con las herramientas convencionales: la ausencia de un interlocutor real que corrija, reactive y exija. Sin supervisión profesional, las rutinas se vuelven repetitivas y pierden progresión técnica; sin variabilidad, el usuario las abandona. Esta dinámica no es anecdótica: los estudios sobre la interacción humano-computador documentan sistemáticamente las altas tasas de deserción en aplicaciones de fitness basadas en pantallas bidimensionales, atribuidas a la monotonía perceptual y a la incapacidad del medio para ofrecer retroalimentación biomecánica o técnica en tiempo real. La consecuencia no es solo la pérdida de motivación: la falta de orientación sobre la distancia, el ritmo y la postura de combate introduce riesgos reales de lesiones menores que ningún video en bucle puede prevenir.

La Realidad Virtual no es simplemente una pantalla más grande o más cercana, es un cambio de paradigma en la relación entre el cuerpo del usuario y el entorno digital. Al posibilitar entornos inmersivos respaldados por seguimiento de seis grados de libertad, elimina la mediación abstracta de los controles convencionales y convierte el movimiento físico real en la interfaz principal. Tomando como punto de partida los principios de diseño de interfaces inmersivas y locomoción espacial trabajados durante el curso, nuestro proyecto aprovecha la plataforma Meta Quest para transformar cualquier habitación en un gimnasio de entrenamiento reactivo y adaptativo. La integración de física tridimensional en tiempo real, retroalimentación háptica diferenciada en los mandos e interfaces completamente diegéticas durante el combate, sin menús flotantes ni interrupciones visuales, permite que la propuesta trascienda el terreno del entretenimiento de arcade para operar como una herramienta genuina de acondicionamiento físico y gestión del estrés, fundamentada en principios sólidos de usabilidad y presencia espacial.

% ============================================================
\section{Requisitos de Diseño y Persona}
\subsection{Requisitos Funcionales}
Los requisitos se derivaron de un análisis de necesidades centrado en el usuario objetivo: jóvenes adultos que buscan ejercicio intensivo e inmersivo sin acceso frecuente a un gimnasio. Cada requisito está redactado como un enunciado único y verificable, con un identificador, y agrupado por categoría. Las justificaciones de diseño se desarrollan en las secciones~\ref{sec:marco} a~\ref{sec:plan}.

\textbf{Interacción y control}\par\nopagebreak
\tablita{|c|L|}{\hd{ID}&\hd{Requisito}}{
RF-01 & El sistema deberá renderizar la vista en primera persona y los guantes virtuales a partir del seguimiento 6DoF del visor y los controladores Meta Quest.\\ \hline
RF-02 & El sistema deberá detectar la colisión entre cada guante del jugador y el oponente, y registrarla como golpe conectado.\\ \hline
RF-03 & El sistema deberá clasificar cada golpe conectado como jab, cross o uppercut según la trayectoria y la velocidad del guante en el instante del impacto.\\ \hline
RF-04 & El sistema deberá reducir la salud del oponente en proporción al tipo de golpe y a la velocidad del impacto.\\ \hline
RF-05 & El sistema deberá detectar un bloqueo cuando ambos guantes del jugador estén en posición de guardia frente al rostro al recibir un ataque, y deberá reducir el daño recibido.\\ \hline
RF-06 & El sistema deberá detectar un esquive cuando la cabeza del jugador se encuentre fuera de la zona de impacto del ataque rival en el instante en que este se ejecuta, y deberá anular el daño de ese ataque.\\
}

\textbf{Oponentes y progresión}\par\nopagebreak
\tablita{|c|L|}{\hd{ID}&\hd{Requisito}}{
RF-07 & El sistema deberá incluir cinco oponentes controlados por IA (Glass Joe, Bald Bull, Soda Popinski, Mr.\ Sandman e Iron Mike Tyson), cada uno con patrón de ataque, velocidad y probabilidad de bloqueo o esquive propios.\\ \hline
RF-08 & El sistema deberá aumentar la dificultad de forma ascendente del Nivel 1 al Nivel 5.\\ \hline
RF-09 & El sistema deberá habilitar el siguiente nivel únicamente cuando el jugador derrote al oponente del nivel actual.\\
}

\textbf{Estructura del combate}\par\nopagebreak
\tablita{|c|L|}{\hd{ID}&\hd{Requisito}}{
RF-10 & El sistema deberá organizar cada combate en 3 rounds de 90 segundos, con señal sonora de campana al inicio y al final de cada round.\\ \hline
RF-11 & El sistema deberá declarar ganador por KO cuando la salud de uno de los contendientes llegue a cero, o por mayor salud restante al terminar el tercer round.\\ \hline
RF-12 & El sistema deberá iniciar una cuenta regresiva del réferi de 10 segundos cuando la salud del jugador llegue a cero, y deberá permitirle levantarse si presiona repetidamente un botón del controlador dentro de ese tiempo.\\ \hline
RF-13 & El sistema deberá declarar derrota si el jugador no supera la cuenta del réferi, y deberá permitirle reintentar el mismo nivel.\\
}

\textbf{Retroalimentación}\par\nopagebreak
\tablita{|c|L|}{\hd{ID}&\hd{Requisito}}{
RF-14 & El sistema deberá comunicar la salud del jugador mediante oscurecimiento progresivo de la visión periférica y un sonido de latido cardíaco cuya frecuencia aumente a medida que la salud disminuye, sin mostrar HUD durante el combate.\\ \hline
RF-15 & El sistema deberá emitir vibración háptica con patrones distintos para golpe conectado, golpe bloqueado y golpe recibido.\\ \hline
RF-16 & El sistema deberá reproducir audio espacial 3D para los golpes, la multitud y la voz del réferi, posicionado según la fuente en el ring.\\
}

\textbf{Menús, resultados y seguridad}\par\nopagebreak
\tablita{|c|L|}{\hd{ID}&\hd{Requisito}}{
RF-17 & El sistema deberá ofrecer un menú principal diegético con las opciones Jugar, Opciones y Salir, seleccionables mediante un puntero láser que sale del control del jugador.\\ \hline
RF-18 & El sistema deberá mostrar, al terminar cada combate, la precisión de golpes, los golpes bloqueados y el tiempo total.\\ \hline
RF-19 & El sistema deberá mostrar el cinturón de campeón al derrotar al oponente del Nivel 5.\\ \hline
RF-20 & El sistema deberá pausar el combate cuando se active el sistema Guardián de Meta Quest y reanudarlo cuando el jugador regrese al área segura.\\ \hline
RF-21 & El sistema deberá registrar por sesión la duración, los combates reintentados, la precisión de golpes y los golpes bloqueados, para su análisis posterior.\\
}

\subsection{Requisitos No Funcionales}
\textbf{Rendimiento y usabilidad}\par\nopagebreak
\tablita{|c|L|}{\hd{ID}&\hd{Requisito}}{
RNF-01 & El sistema deberá mantener un mínimo de 72 FPS en Meta Quest 2 durante el combate, sin caídas por debajo de 72 FPS en más del 5\,\% de los fotogramas de una sesión.\\ \hline
RNF-02 & El sistema deberá iniciar la respuesta háptica en un máximo de 15 ms después de detectada la colisión.\\ \hline
RNF-03 & Al menos el 80\,\% de los usuarios novatos deberán adoptar la postura de guardia y ejecutar un golpe válido en un máximo de 2 minutos de juego libre, sin instrucciones externas.\\ \hline
RNF-04 & El sistema no deberá emplear teletransporte ni locomoción por joystick; el desplazamiento será exclusivamente físico a escala 1:1.\\ \hline
RNF-05 & Los usuarios deberán calificar la ausencia de cinetosis con un promedio mínimo de 4 sobre 5 en un cuestionario Likert posterior a la prueba.\\
}

\textbf{Seguridad, fiabilidad, compatibilidad y mantenibilidad}\par\nopagebreak
\tablita{|c|L|}{\hd{ID}&\hd{Requisito}}{
RNF-06 & El sistema deberá operar dentro de un área despejada mínima de 1,5 $\times$ 1,5 m y respetar los límites definidos por el sistema Guardián.\\ \hline
RNF-07 & El sistema deberá registrar al menos el 95\,\% de los impactos reales de los guantes sobre el oponente, incluso en golpes de alta velocidad.\\ \hline
RNF-08 & El sistema deberá soportar 30 minutos de sesión continua sin cierres inesperados ni pérdida de seguimiento atribuible a la aplicación.\\ \hline
RNF-09 & El sistema deberá ejecutarse de forma autónoma (standalone) en Meta Quest 2 y Meta Quest 3, sin requerir una PC conectada.\\ \hline
RNF-10 & Los parámetros de cada oponente (velocidad, tipo de ataque, probabilidad de bloqueo y esquive) deberán estar definidos como datos configurables, de modo que se puedan ajustar sin modificar el código del sistema de combate.\\
}

\subsection{Especificación de Persona}
\tablita{|>{\bfseries}p{4.2cm}|L|}{\hd{Atributo}&\hd{Descripción}}{
Nombre & Santiago Miguel Silva Reyes\\ \hline
Edad & 20 años\\ \hline
Ocupación & Estudiante de Ingeniería\\ \hline
Motivación principal & Mantenerse activo y canalizar el estrés académico sin salir de su habitación\\ \hline
Disponibilidad & Limitada (no puede acudir regularmente a un gimnasio convencional)\\ \hline
Intereses & Tecnología, videojuegos interactivos, deportes de contacto\\ \hline
Frustración actual & Monotonía de rutinas en video y falta de sensación de reto físico real en pantallas planas\\ \hline
Dispositivo & Meta Quest 2\\ \hline
Expectativa & Experiencia inmersiva de alta intensidad física con control intuitivo y sin menús complejos\\
}

\subsection{Narrativa y Objetivo de la Experiencia}
La narrativa sitúa al usuario en el papel de un peleador novato que ingresa al cuadrilátero de un gimnasio underground. Con la meta de obtener el cinturón de campeón regional, el protagonista debe medirse consecutivamente contra los cinco peleadores más destacados de la liga local, cada uno presentando una técnica de combate distintiva.

El objetivo fundamental es transmitir la adrenalina, tensión y presencia espacial propias de un combate de boxeo real, garantizando un esquema de control completamente orgánico libre de menús flotantes que interrumpan el flujo de la actividad física durante el combate.

\subsection{Arco Narrativo}
\label{sec:arco}
El arco sigue la pirámide de Freytag (exposición, acción ascendente, clímax, acción descendente y desenlace), con el detonante señalado como un punto propio entre la exposición y la acción ascendente. La tabla~\ref{tab:arco} relaciona cada etapa con lo que ocurre en la experiencia y con el nivel del juego en que sucede.

\begin{table}[H]
\centering
\caption{Arco narrativo de la experiencia según la pirámide de Freytag.}
\label{tab:arco}
\tablita{|>{\bfseries}p{3.2cm}|L|p{2.6cm}|}{\hd{Etapa}&\hd{En la experiencia}&\hd{Nivel}}{
Exposición & El jugador entra a la arena del gimnasio underground (el menú principal está dentro del ring), ve a los rivales y aprende el control con los guantes: agacharse esquiva, cubrir la cara bloquea y extender el brazo golpea. & Menú principal\\ \hline
Detonante (incidente incitador) & Suena la campana del primer round contra Glass Joe, oponente introductorio de velocidad moderada y patrones simples: empieza el camino al cinturón regional. & Nivel 1\\ \hline
Acción ascendente & Cada rival sube la presión y cada victoria desbloquea al siguiente: Bald Bull con ataques directos de fuerza, Soda Popinski con combinaciones rápidas de ganchos y Mr.\ Sandman con técnica de alta velocidad y esquive activo. Las caídas y la cuenta del réferi son los puntos de tensión. & Niveles 2, 3 y 4\\ \hline
Clímax & Enfrentamiento definitivo contra Iron Mike Tyson, el campeón invicto: máxima intensidad, el menor descuido resulta en KO. & Nivel 5\\ \hline
Acción descendente & La caída final de Tyson, la cuenta del réferi que ya no se supera y la ovación de la multitud. & Nivel 5\\ \hline
Desenlace & Se despliega el cinturón de campeón y la hoja de resumen estadístico del desempeño (precisión de golpes, golpes bloqueados, tiempo total). & Resultados\\
}
\end{table}

La figura~\ref{fig:freytag} dibuja el mismo arco como la pirámide de Freytag vista en clase: la tensión sube con cada nivel hasta el combate contra Tyson y baja con su caída hasta el desenlace en la escena de resultados.

\begin{figure}[H]
  \centering
  \begin{tikzpicture}[x=1cm,y=1cm,font=\small]
    % ejes
    \draw[->,gray] (0,0) -- (15.2,0) node[below left,black]{tiempo de la experiencia};
    \draw[->,gray] (0,0) -- (0,5.2) node[above right,black]{tensión};
    % pirámide
    \draw[very thick,vino] (0.2,0.6) -- (2.6,0.6) -- (3.4,1.2) -- (9.6,4.4) -- (12.2,1.6) -- (14.8,1.0);
    \fill[dorado] (3.4,1.2) circle (2.4pt);
    \fill[rojo] (9.6,4.4) circle (3pt);
    % niveles sobre la acción ascendente
    \foreach \x/\y/\n in {4.9/2.0/2, 6.4/2.75/3, 7.9/3.5/4} {\fill[vino] (\x,\y) circle (1.8pt); \node[above left=-1pt] at (\x,\y) {\scriptsize N\n};}
    % etiquetas
    \node[align=center,below] at (1.4,0.55) {\textbf{Exposición}\\Menú en la arena};
    \node[align=center,above left] at (3.5,1.25) {\textbf{Detonante}\\Campana, nivel 1};
    \node[align=center,rotate=27] at (7.0,2.45) {\textbf{Acción ascendente} (niveles 2 a 4)};
    \node[align=center,above] at (9.6,4.5) {\textbf{Clímax}: Iron Mike Tyson (nivel 5)};
    \node[align=center,rotate=-47] at (11.6,3.3) {\textbf{Acción descendente}\\[-1pt]\scriptsize caída final, cuenta y ovación};
    \node[align=center,above] at (14.1,1.15) {\textbf{Desenlace}};
    \node[align=center,below] at (14.1,0.95) {\scriptsize cinturón y estadísticas};
  \end{tikzpicture}
  \caption{Arco narrativo de la experiencia sobre la pirámide de Freytag (N2 a N4: niveles 2 a 4).}
  \label{fig:freytag}
\end{figure}

\begin{figure}[H]
  \centering
  \begin{subfigure}[b]{0.38\linewidth}
    \centering\includegraphics[height=6.2cm]{glass_joe.png}
    \caption{Glass Joe, Nivel 1 (detonante).}
  \end{subfigure}\hfill
  \begin{subfigure}[b]{0.38\linewidth}
    \centering\includegraphics[height=6.2cm]{iron_mike_tyson.png}
    \caption{Iron Mike Tyson, Nivel 5 (clímax).}
  \end{subfigure}
  \caption{Los rivales que abren y cierran el arco narrativo.}
  \label{fig:rivales-arco}
\end{figure}

\subsection{Métrica HEART}
Para evaluar la usabilidad y adopción de la experiencia se adoptó el framework HEART, marco aprendido durante el curso y reconocido en la industria del diseño centrado en el usuario. Las pruebas con usuarios se hacen en la semana 14 del cronograma (sección~\ref{sec:crono}) con el protocolo de la sección~\ref{sec:heart}. Desde el Entregable~4 el juego guarda por sí mismo los datos de \emph{Engagement} y \emph{Task Success} (sección~\ref{sec:registro}), y desde el Entregable~5 también el tiempo hasta la primera guardia y el primer golpe, que respalda la observación de \emph{Adoption}. La fila de \emph{Retention} se adaptó: el número de sesiones en la segunda y la tercera semana no se puede medir en la semana 14, porque los participantes no tienen visor propio y la entrega es esa misma semana.

\tablita{|>{\bfseries}p{2.4cm}|L|L|}{\hd{Dimensión HEART}&\hd{Indicador de medición}&\hd{Instrumento}}{
Happiness & Satisfacción inmersiva percibida y ausencia de cinetosis & Cuestionario post-prueba (escala Likert 1--5)\\ \hline
Engagement & Duración media de sesión y número de combates reintentados & Registro de sesión del juego (JSON y CSV, sección~\ref{sec:registro})\\ \hline
Adoption & Tiempo para comprender la postura de guardia sin tutorial & Observación directa (cronómetro) y registro del juego (tiempo hasta la guardia y el primer golpe)\\ \hline
Retention & Intención de volver a jugar y continuación voluntaria al ofrecerle parar (en lugar de las sesiones en segunda y tercera semana) & Cuestionario post-prueba y observación\\ \hline
Task Success & \% de golpes bloqueados con éxito, precisión de impacto, superación de la cuenta del réferi & Métricas registradas por el sistema (\cs{RegistroSesion})\\
}

% ============================================================
\section{Marco de Diseño}
\label{sec:marco}
\subsection{Fundamentos del Framework de Interacción}
El framework de interacción y diseño visual se fundamenta en dos principios centrales estudiados en el curso: la interacción corporizada (\emph{embodied interaction}) y la manipulación directa tridimensional. En lugar de proyectar interfaces abstractas bidimensionales en el campo de visión del usuario, la arquitectura de diseño integra toda la información crítica de forma diegética en el entorno virtual.

Tal como expone el concepto de conformidad sensorial trabajado en clase, el diseño busca minimizar la brecha entre la acción física del usuario y la respuesta del sistema virtual. Al mapear directamente la posición de las manos a los guantes virtuales a través del seguimiento 6DoF, se eliminan los intermediarios de entrada abstractos: agacharse equivale a esquivar, cubrir la cara equivale a bloquear y extender el brazo equivale a golpear. No existe capa de traducción simbólica.

\subsection{Principios de Diseño Aplicados}
\begin{itemize}
  \item \textbf{Reducción de carga cognitiva:} toda la información de estado (salud, cuenta del réferi) se comunica mediante señales visuales y auditivas diegéticas integradas en el entorno. No se utilizan HUD flotantes ni menús bidimensionales durante el combate.
  \item \textbf{Retroalimentación multimodal cerrada:} el bucle de retroalimentación sensorial combina el oscurecimiento periférico de la pantalla, la respuesta sonora del pulso cardíaco y la vibración háptica diferenciada de los mandos para comunicar el estado del jugador sin desviar su atención del oponente.
  \item \textbf{Presencia e inmersión:} la perspectiva visual en primera persona a escala 1:1, el sonido espacial tridimensional y la respuesta háptica coordinada construyen un nivel de presencia elevado, consistente con los modelos de presencia de Slater y Wilbur revisados en el curso.
  \item \textbf{Seguridad física:} la locomoción estrictamente física con respeto al sistema Guardián aplica el principio de diseño responsable en entornos de RV de espacio completo.
  \item \textbf{Curva de aprendizaje mínima:} el control se aprende mediante la acción directa. La ausencia de tutoriales extensos responde al principio de \emph{affordance} natural: el usuario sabe que debe golpear porque ve guantes en sus manos.
\end{itemize}

\subsection{Implementación Técnica del Framework}
La arquitectura de interacción se implementa sobre el Meta XR SDK (paquete \cs{com.meta.xr.sdk.all}) con OpenXR como sistema de entrada. Estos son los componentes que sostienen la interacción.

\textbf{OVR Camera Rig.} Provee el sistema de cámara principal anclado a la cabeza del usuario con seguimiento 6DoF en tiempo real. Sincroniza la perspectiva visual en primera persona con los movimientos físicos de cabeza y torso, siendo el componente central de la presencia espacial. Cada movimiento de la cabeza del jugador se traduce inmediatamente en la orientación de la cámara virtual dentro del ring.

\textbf{Guantes sobre los anclajes de las manos.} Los guantes virtuales son hijos de los anclajes de mano del \cs{OVRCameraRig}, de modo que la posición y la rotación de los mandos Meta Quest se reflejan 1:1 en los guantes y, incluso en movimientos rápidos de puño, no hay separación entre la mano virtual y la física. La detección de los golpes la realizan los colliders de los guantes y del oponente (RF-02 y RF-03), no un componente de agarre.

\textbf{Interacción con los menús.} El menú diegético usa el sistema de interfaz de Meta Core (\cs{OVRInputModule} y \cs{OVRRaycaster}) con un puntero láser que sale del control derecho y se acciona con el gatillo.

\textbf{Hand Grab Recorder propio.} El informe planteó registrar y reproducir secuencias de movimiento de mano reales para la animación procedural de los bloqueos de la IA, superando las limitaciones de las animaciones esqueléticas estáticas. Se implementó como un grabador propio (\cs{GrabadoraGuardia}, sección~\ref{sec:grabadora}): graba con el visor puesto los movimientos de ambas manos respecto a la cabeza, y \cs{DefensaOponente} los reproduce mediante cinemática inversa en los rivales.

\subsection{Bocetos de Interfaz y Flujo de Experiencia}
Se elaboraron dos bocetos en primera persona que muestran las dos situaciones centrales de la experiencia: el combate y el estado de knockout. Ambos ilustran la ausencia total de elementos de HUD superpuestos durante el combate; toda la información se transmite mediante el entorno, el sonido y la vibración.

\begin{figure}[H]
  \centering
  \includegraphics[width=0.92\linewidth]{boceto_combate.png}
  \caption{Boceto de la escena de combate (\cs{Scene_Ring_Level1}, 1P vs.\ IA, primera persona, sin HUD).}
  \label{fig:boceto-combate}
\end{figure}

\tablita{|c|>{\bfseries}p{3.4cm}|L|}{\hd{N.º}&\hd{Elemento}&\hd{Descripción}}{
1 & Esquive activo & Mueve cabeza y torso para evadir los golpes.\\ \hline
2 & Oponente IA & Glass Joe, nivel 1, controlado por la IA. Cada rival tiene un patrón único.\\ \hline
3 & Tus guantes & Siguen tus manos en 1:1 (seguimiento 6DoF).\\ \hline
4 & Salud & Bordes oscuros: a menos salud, más se oscurece la vista.\\ \hline
5 & Pulso cardíaco & Se escucha y refuerza el estado de salud y resistencia.\\ \hline
6 & Réferi & Cuenta ante un KO.\\ \hline
7 & Vibración & Los mandos vibran distinto: golpe conectado o bloqueo.\\
}

\begin{figure}[H]
  \centering
  \includegraphics[width=0.92\linewidth]{boceto_knockout.png}
  \caption{Boceto del estado de knockout (cuenta del réferi, 1P vs.\ IA, sin HUD).}
  \label{fig:boceto-ko}
\end{figure}

\tablita{|c|>{\bfseries}p{3.4cm}|L|}{\hd{N.º}&\hd{Elemento}&\hd{Descripción}}{
1 & Salud crítica & La vista casi se cierra: los bordes quedan muy oscuros.\\ \hline
2 & Pulso cardíaco & Se escucha acelerado mientras estás en la lona.\\ \hline
3 & Réferi & Hace la cuenta regresiva tras el KO.\\ \hline
4 & Button mashing & Pulsa rápido los botones del mando para levantarte.\\
}

Superar la cuenta del réferi permite continuar el combate. En ambos bocetos la locomoción es 1:1 y respeta el sistema Guardián de Meta Quest, sin barras ni menús flotantes durante el combate. Estos bocetos son la referencia del Entregable~6, donde se compara lo diseñado en este hito con lo efectivamente implementado. La sección~\ref{sec:bocetos} hace esa comparación elemento por elemento, con capturas del juego; las demás escenas están en la sección~\ref{sec:estructura}.

% ============================================================
\section{Prototipo Digital}
\subsection{Descripción General del Prototipo}
El prototipo digital se desarrolla en el motor de videojuegos Unity 2022 LTS utilizando C\# como lenguaje principal para la gestión de estados del combate y la inteligencia artificial reactiva. El destino de compilación es Android (Meta Quest) y la interacción inmersiva usa OpenXR con el Meta XR SDK, que garantiza la integración del seguimiento de seis grados de libertad.

El entorno visual usa el pipeline de renderizado integrado (\emph{Built-in}) con materiales Standard y un escenario de gimnasio de boxeo (arena de MarpaStudio), con el objetivo de sostener un mínimo de 72 FPS (RNF-01) en el hardware objetivo.

El desarrollo es incremental: cada entregable del cronograma (sección~\ref{sec:crono}) añade capacidades verificables, desde un proyecto compilando en Meta Quest con el ring a escala 1:1 (15\,\%), pasando por los primeros golpes con háptica contra Glass Joe (30\,\%), combates completos contra dos rivales con salud diegética (50\,\%) y cuatro rivales con ciclo de KO (70\,\%), hasta la versión final con los cinco niveles y las pruebas de usabilidad (100\,\%, este entregable).

% ============================================================
\section{Iniciando el Desarrollo: Escenas y Assets en Unity}
\subsection{Estructura del Proyecto en Unity}
\label{sec:estructura}
El proyecto está organizado en Unity 2022 LTS con las siete escenas del diseño original incluidas en los \emph{Build Settings}; las dos últimas se incorporaron en el Entregable~5.

\begin{itemize}
  \item \cs{Main_Scene}: menú diegético tridimensional dentro de la misma arena del ring, con música de fondo en bucle (botón Jugar con selección de nivel, y Opciones de volumen y háptica). Los boxeadores que se ven en ella son decoración.
  \item \cs{Scene_Ring_Level1}: cuadrilátero con el oponente Glass Joe.
  \item \cs{Scene_Ring_Level2}: cuadrilátero con Bald Bull.
  \item \cs{Scene_Ring_Level3}: cuadrilátero con Soda Popinski.
  \item \cs{Scene_Ring_Level4}: cuadrilátero con Mr.\ Sandman.
  \item \cs{Scene_Ring_Level5}: cuadrilátero con Iron Mike Tyson (Entregable~5, sección~\ref{sec:tyson}).
  \item \cs{Scene_Results}: hoja de resultados tras cada combate y despliegue del cinturón de campeón (Entregable~5, sección~\ref{sec:resultados}).
\end{itemize}

La figura~\ref{fig:escenas} muestra las escenas implementadas, capturadas desde los ojos del jugador.

\begin{figure}[H]
  \centering
  \begin{subfigure}[t]{0.48\linewidth}
    \includegraphics[width=\linewidth]{01_menu_principal.png}
    \caption{Menú principal diegético dentro de la arena.}
  \end{subfigure}\hfill
  \begin{subfigure}[t]{0.48\linewidth}
    \includegraphics[width=\linewidth]{02_nivel1_glass_joe.png}
    \caption{Nivel 1: Glass Joe.}
  \end{subfigure}\\[6pt]
  \begin{subfigure}[t]{0.48\linewidth}
    \includegraphics[width=\linewidth]{03_nivel2_bald_bull.png}
    \caption{Nivel 2: Bald Bull.}
  \end{subfigure}\hfill
  \begin{subfigure}[t]{0.48\linewidth}
    \includegraphics[width=\linewidth]{04_nivel3_soda_popinski.png}
    \caption{Nivel 3: Soda Popinski.}
  \end{subfigure}\\[6pt]
  \begin{subfigure}[t]{0.48\linewidth}
    \includegraphics[width=\linewidth]{05_nivel4_mr_sandman.png}
    \caption{Nivel 4: Mr.\ Sandman.}
  \end{subfigure}\hfill
  \begin{subfigure}[t]{0.48\linewidth}
    \includegraphics[width=\linewidth]{11_nivel5_mike_tyson.png}
    \caption{Nivel 5: Iron Mike Tyson.}
  \end{subfigure}\\[6pt]
  \begin{subfigure}[t]{0.48\linewidth}
    \includegraphics[width=\linewidth]{13_resultados_victoria.png}
    \caption{Hoja de resultados tras una victoria.}
  \end{subfigure}\hfill
  \begin{subfigure}[t]{0.48\linewidth}
    \includegraphics[width=\linewidth]{15_resultados_campeon.png}
    \caption{Cinturón de campeón tras ganar el nivel 5.}
  \end{subfigure}
  \caption{Escenas implementadas del proyecto (capturas en Play, desde los ojos del jugador).}
  \label{fig:escenas}
\end{figure}

\subsection{Tabla de Assets}
\tablita{|>{\raggedright\arraybackslash}p{2.4cm}|>{\raggedright\arraybackslash}p{4.6cm}|>{\raggedright\arraybackslash}p{2.6cm}|L|}{\hd{Elemento del juego}&\hd{Asset / Modelo 3D}&\hd{Origen}&\hd{Función en la experiencia}}{
Escenario principal & Boxing Arena Game Level / Boxing Ring Gym & Sketchfab y MarpaStudio & Cuadrilátero 3D con iluminación focalizada y cuerdas delimitadoras\\ \hline
Oponente Nivel 1 & Glass Joe 3D Model (rigged) & Sketchfab, rig de Mixamo & Rival introductorio de baja velocidad y patrones simples\\ \hline
Oponente Nivel 2 & Bald Bull 3D Model (rigged) & Sketchfab, rig de Mixamo & Rival de fuerza intermedia enfocado en ataques directos\\ \hline
Oponente Nivel 3 & Soda Popinski 3D Model (rigged) & Sketchfab, rig de Mixamo & Rival ágil con combinaciones rápidas de ganchos\\ \hline
Oponente Nivel 4 & Mr.\ Sandman 3D Model & Sketchfab, riggeado por el equipo en Blender & Rival técnico de alta velocidad y esquive activo\\ \hline
Oponente Nivel 5 & Iron Mike Tyson 3D Model & Sketchfab, riggeado por el equipo & Rival final de máxima velocidad, agresividad y precisión (malla reducida a 15\,000 triángulos)\\ \hline
Réferi & Boxing Referee 3D Model (rigged) & Sketchfab & Árbitro que realiza la cuenta regresiva ante un KO\\ \hline
Guantes del jugador & Boxing Gloves (low poly) & Sketchfab & Representación visual de las manos del jugador\\ \hline
Cinturón de campeón & Modelo propio del cinturón, generado por \cs{Tools/crear_cinturon.py} & Propio: generado con Blender por script, sin licencia de terceros & Objeto de victoria al finalizar el Nivel 5 (Entregable~5)\\
}

\subsection{Configuración de los Componentes OVR en Unity}
\textbf{OVR Camera Rig.} El prefab \cs{OVRCameraRig} (Meta XR SDK) está configurado con \emph{Tracking Origin Mode} en \emph{Floor Level} para que la posición de los guantes virtuales corresponda a la posición real de las manos del jugador respecto al suelo. El componente \cs{OVRManager} tiene habilitado el soporte de controladores para Meta Quest 2 y Meta Quest 3.

\textbf{Guantes.} Cada guante virtual es hijo del anclaje de la mano correspondiente del \cs{OVRCameraRig}, bajo un objeto \cs{Puno} que lleva el \emph{Rigidbody} cinemático y el collider. Esto sigue la posición y rotación del control 1:1 y el guante no se suelta durante el combate. El guante no interviene en la detección de impactos más allá de su collider: el contacto se resuelve con scripts de colisión sobre los guantes y el oponente.

\textbf{Hand Grab Recorder.} Se utiliza en dos contextos: durante el desarrollo, para capturar y validar patrones de movimiento de guardia, y en la animación de bloqueos de los oponentes, donde las secuencias de mano grabadas se reproducen proceduralmente para dotar a los rivales de transiciones de bloqueo más naturales que las de las animaciones esqueléticas estáticas. Se implementó en el Entregable~4 (sección~\ref{sec:grabadora}).

\subsection{Implementación del Entregable 2 (Semana 11)}
El Entregable~2 corresponde al 30\,\% de avance acumulado y quedó completado. Su resultado verificable es que el jugador golpea a Glass Joe y siente la vibración en el control, y que los paneles de juego y de opciones se abren desde el menú principal. A continuación se detalla cómo se implementó cada actividad en Unity, con los scripts en C\# que la resuelven (carpeta \cs{Assets/Scripts}).

\subsubsection*{Menú principal: panel de juego y panel de opciones}
El menú principal es un \emph{Canvas} en \emph{World Space} ubicado dentro de la arena (escena \cs{Main_Scene}), con música de fondo en bucle. El jugador apunta con un láser que sale del control derecho y selecciona con el gatillo. La interacción usa \cs{OVRInputModule} y \cs{OVRRaycaster} de Meta XR Core SDK, y el script \cs{LaserCursor} dibuja el rayo y un punto en el lugar donde apunta. El script \cs{MainMenu} gestiona los paneles:
\begin{itemize}
  \item \textbf{Panel de juego:} el botón JUGAR abre la selección de nivel, cuyos botones cargan cada combate. El nivel 1 está habilitado al inicio de la progresión y cada nivel siguiente se abre al derrotar al anterior (RF-09).
  \item \textbf{Panel de opciones:} el botón OPCIONES cambia al panel de opciones, con los controles de volumen y de háptica, y ATRÁS regresa al panel principal. Los ajustes se guardan con \cs{PlayerPrefs}: \cs{"volumen"} se aplica a \cs{AudioListener.volume} al iniciar y \cs{"haptica"} escala la intensidad de todas las vibraciones de los guantes.
  \item \textbf{Salir:} el botón SALIR cierra la aplicación.
\end{itemize}
Dentro del ring, el script \cs{PauseMenu} abre el mismo Canvas como menú de pausa con el botón Menú del control izquierdo. Detiene el combate (\cs{Time.timeScale = 0}) y ofrece REANUDAR y VOLVER AL MENÚ PRINCIPAL.

\subsubsection*{Guantes del jugador}
Los guantes provienen del modelo \cs{Boxing gloves.glb}, importado con glTFast. Una herramienta de editor propia (\emph{Tools > Boxeo > Poner guantes en las manos}, script \cs{GuantesBuilder}) separa el modelo en dos mallas, \cs{Guante_Izquierdo} y \cs{Guante_Derecho}, y coloca cada una en la mano correspondiente del \cs{OVRCameraRig}. Están giradas para que el pulgar quede hacia arriba al sostener el control, como en la guardia real.

\subsubsection*{Detección de colisiones y clasificación de jab, cross y uppercut}
Cada guante lleva el script \cs{Puno}, con un \emph{Rigidbody} cinemático y un collider con detección continua, para que los golpes rápidos no atraviesen al oponente (RNF-07). La velocidad de la mano se mide respecto al espacio de tracking del jugador, así que lo que cuenta es el movimiento real del brazo. Un contacto se registra como golpe conectado (RF-02) solo si la mano avanza hacia el oponente a 0,5 m/s o más. Después de cada golpe, ese puño espera 0,25 s antes de poder conectar otro, para que un mismo impacto no se cuente dos veces.

El oponente tiene dos zonas de impacto (script \cs{ZonaGolpe}): Cabeza y Torso. Son colliders de tipo \emph{trigger} emparentados a los huesos de la cabeza y de la columna, así que acompañan cada animación. La herramienta \emph{Tools > Boxeo > Configurar oponente seleccionado} las crea a partir de los vértices del modelo asociados a esos huesos.

En el instante del impacto, \cs{Golpe.Clasificar} determina el tipo de golpe (RF-03) a partir de la dirección de la mano respecto a hacia dónde mira el jugador, con guardia ortodoxa (diestra):

\begin{table}[H]
\centering
\caption{Clasificación de los golpes del jugador.}
\tablita{|>{\bfseries}p{2.6cm}|L|p{3.4cm}|}{\hd{Tipo de golpe}&\hd{Condición de detección}&\hd{Mano}}{
Uppercut & La mano sube con menos de 50° respecto a la vertical. & Cualquiera\\ \hline
Jab & Golpe recto: desviación lateral menor a 40° respecto al frente. & Izquierda (adelantada)\\ \hline
Cross & Golpe recto: desviación lateral menor a 40° respecto al frente. & Derecha (atrasada)\\ \hline
Gancho (lateral) & Cualquier otra trayectoria, de lado. & Cualquiera\\
}
\end{table}

Cada golpe se registra en la consola de Unity con su tipo, mano, zona, velocidad y el conteo acumulado, lo que permite validar la detección durante las pruebas.

\subsubsection*{Glass Joe en el ring}
Glass Joe es un modelo con esqueleto humanoide de Mixamo (1,70 m) ubicado en el centro del ring en la escena \cs{Scene_Ring_Level1}. Su \emph{Animator Controller} reúne la guardia (\cs{Idle}), las reacciones a los golpes, las caídas y sus propios ataques. El script \cs{Oponente} lleva su salud y lo hace caer al cruzar umbrales: con tres caídas por derrota, las dos primeras son knockdowns y la tercera es la derrota. Mientras está en la lona no recibe golpes. El script \cs{ReaccionAnimada} elige la animación según el golpe y la zona:

\begin{table}[H]
\centering
\caption{Reacciones animadas de Glass Joe.}
\tablita{|L|L|}{\hd{Golpe recibido}&\hd{Reacción de Glass Joe}}{
Jab, cross o gancho a la cabeza & \cs{Jab} (la cabeza va hacia atrás)\\ \hline
Uppercut a la cabeza & \cs{Uppercut_Face}\\ \hline
Cualquier golpe al torso & \cs{Uppercut_Stomach} (se dobla)\\ \hline
Salud en un umbral (1.ª y 2.ª caída) & \cs{Knockdown} y regreso a la guardia\\ \hline
Salud en cero (3.ª caída) & \cs{Defeat} (queda derrotado)\\
}
\end{table}

Además, el script \cs{AtaqueOponente} hace que Glass Joe gire hacia el jugador y lance un ataque cada 4 a 6 s desde la guardia: jabs y, cada cuatro ataques, un uppercut. Si el jugador lo presiona, cada golpe recibido reinicia la espera y corta el ataque en curso. Los valores del personaje (salud, caídas, tiempos de ataque, sonidos) están en un asset \cs{PerfilOponente} (\cs{Assets/_Project/Data/Oponentes/GlassJoe.asset}) y se ajustan sin modificar el código, en línea con RNF-10.

\subsubsection*{Háptica básica}
Al conectar un golpe, el control de esa mano vibra. Cada tipo de golpe tiene su propio patrón, de modo que el jugador distingue lo que conectó sin mirar (base de RF-15):

\begin{table}[H]
\centering
\caption{Vibración por tipo de golpe conectado.}
\tablita{|>{\bfseries}p{2.2cm}|L|p{3.6cm}|}{\hd{Golpe}&\hd{Patrón de vibración}&\hd{Sensación}}{
Jab & Un pulso de amplitud 0,7 durante 50 ms. & Seco y corto\\ \hline
Cross & Un pulso de amplitud 1,0 durante 140 ms. & Largo y fuerte\\ \hline
Uppercut & Pulso de 0,5 (50 ms), pausa de 40 ms y pulso de 1,0 (100 ms). & Doble, ascendente\\ \hline
Gancho & Un pulso de amplitud 0,35 durante 60 ms. & Suave\\
}
\end{table}

La intensidad crece con la rapidez del puño, desde un mínimo de 40\,\% hasta el máximo a 2,5 m/s, y se multiplica por el ajuste de háptica del panel de opciones. La vibración se detiene aunque el juego se pause justo después del golpe.

\subsubsection*{Verificación y resumen de cumplimiento}
La escena se probó en el editor con Meta XR Simulator y en el visor Meta Quest mediante una compilación para Android (OpenXR). Se verificó el recorrido completo del resultado esperado: desde el menú principal se abren el panel de opciones y el panel de juego; al pulsar JUGAR se entra al ring frente a Glass Joe; cada jab, cross y uppercut que conecta produce la reacción animada del oponente y una vibración propia en el control.

\begin{table}[H]
\centering
\caption{Cumplimiento del Entregable 2.}
\tablita{|L|L|c|}{\hd{Actividad planificada}&\hd{Implementación}&\hd{Estado}}{
Panel de juego (selección de nivel) & Botón JUGAR del menú principal (\cs{MainMenu}). & \ok\\ \hline
Panel de opciones & Panel de opciones con volumen y háptica guardados en \cs{PlayerPrefs}; botón ATRÁS. & \ok\\ \hline
Guantes en las manos & \cs{Guante_Izquierdo} y \cs{Guante_Derecho} adheridos a cada mano del \cs{OVRCameraRig} (\cs{GuantesBuilder}, \cs{Puno}). & \ok\\ \hline
Detección de colisiones para jab, cross y uppercut & \cs{Puno}, \cs{ZonaGolpe} (cabeza y torso) y \cs{Golpe.Clasificar}. & \ok\\ \hline
Glass Joe en el ring & Modelo animado con reacciones, knockdowns, derrota y ataques (\cs{Oponente}, \cs{ReaccionAnimada}, \cs{AtaqueOponente}). & \ok\\ \hline
Háptica básica & Vibración por tipo de golpe, escalada por rapidez y por el ajuste de opciones (\cs{Puno}). & \ok\\ \hline
Resultado verificable & El jugador golpea a Glass Joe y siente la vibración; ambos paneles abren desde el menú. & Cumplido\\
}
\end{table}

% ============================================================
\section{Plan de Diseño}
\label{sec:plan}
\subsection{¿Qué pueden lograr los usuarios?}
La aplicación permite ejecutar un entrenamiento técnico reactivo en tiempo real sin salir del hogar. Específicamente, el jugador puede:
\begin{itemize}
  \item Combatir aplicando estrategias de golpeo directo (jab), cruzado (cross), uppercut y ganchos, con retroalimentación háptica diferenciada.
  \item Esquivar activamente los ataques del rival desplazando el torso y la cabeza en el espacio físico.
  \item Enfrentar cinco estilos de boxeo inspirados en luchadores icónicos, cada uno con un patrón de combate único.
  \item Superar niveles de dificultad con velocidad e intensidad incremental.
  \item Luchar por levantarse de la lona ante la cuenta del réferi mediante interacción física de \emph{button mashing}.
  \item Culminar la progresión obteniendo el cinturón de campeonato regional y consultando sus estadísticas de combate.
\end{itemize}

\subsection{Límites Técnicos y de la Experiencia}
\begin{itemize}
  \item \textbf{Fatiga física rápida:} el constante esquive y golpeo en sesiones prolongadas genera fatiga muscular en brazos y core. El diseño limita cada combate a 3 rounds de 90 segundos para gestionar la intensidad.
  \item \textbf{Espacio físico requerido:} se requiere un área despejada mínima de 1,5 $\times$ 1,5 metros respaldada por el sistema Guardián de Meta Quest.
  \item \textbf{Ausencia de seguimiento de extremidades inferiores:} la evasión está restringida al movimiento de torso y cabeza; no se detecta el movimiento de piernas ni el \emph{footwork}.
  \item \textbf{Dependencia del hardware:} la precisión del sistema depende directamente de la tasa de refresco del visor y del correcto seguimiento de los controladores en condiciones de iluminación adecuadas.
\end{itemize}

\subsection{Profundidad de la Presencia}
El proyecto alcanza un nivel de inmersión elevado (presencia de lugar y plausibilidad, según el modelo de Slater) al sincronizar la perspectiva visual en primera persona a escala 1:1 con la propagación del sonido espacial en tres dimensiones y la respuesta háptica coordinada en los controladores. La combinación de estos tres canales sensoriales (visual, auditivo y táctil) construye el bucle de retroalimentación cerrado que caracteriza las experiencias de RV de alta presencia documentadas en la literatura del curso.

% ============================================================
\section{Cronograma de Trabajo}
\label{sec:crono}
El proyecto se ejecuta entre las semanas 10 y 15 mediante seis entregables. Los porcentajes indican el avance acumulado del proyecto en cada entregable.

\subsection{Entregables y avance acumulado}
\begin{table}[H]
\centering
\caption{Entregables del proyecto.}
\label{tab:entregables}
\tablita{|>{\raggedright\arraybackslash}p{1.9cm}|c|c|L|>{\raggedright\arraybackslash}p{3.3cm}|}{\hd{Entregable}&\hd{Semana}&\hd{Avance}&\hd{Actividades principales}&\hd{Resultado verificable}}{
Entregable 1 & 10 & 15\,\% & Proyecto en Unity con Meta XR Core SDK. OVR Camera Rig en Floor Level y compilación de prueba en Meta Quest. Menú principal con los botones Jugar y Opciones. Ring a escala 1:1. & Proyecto compilando en el visor, menú navegable y ring a escala 1:1.\\ \hline
Entregable 2 & 11 & 30\,\% & Panel de juego (selección de nivel) y panel de opciones. Guantes en las manos. Detección de colisiones para jab, cross y uppercut. Glass Joe en el ring y háptica básica. & El jugador golpea a Glass Joe y siente la vibración; ambos paneles abren desde el menú.\\ \hline
Entregable 3 & 12 & 50\,\% & IA reactiva de los niveles 1 y 2 (Glass Joe y Bald Bull) con ataque, bloqueo y velocidad. Esquive con cabeza y torso. Salud diegética con viñeta y pulso cardíaco. & Combates completos contra dos rivales, con la salud percibida sin HUD.\\ \hline
Entregable 4 & 13 & 70\,\% & IA de los niveles 3 y 4 (Soda Popinski y Mr.\ Sandman). Réferi con cuenta regresiva y button mashing tras el KO. Rondas de 3 $\times$ 90 s. Hand Grab Recorder en los bloqueos y háptica diferenciada. & Cuatro rivales jugables y ciclo de KO funcionando.\\ \hline
\rowcolor{crema}\textbf{Entregable 5} & \textbf{14} & \textbf{100\,\%} & IA del nivel 5 (Iron Mike Tyson). Escena de resultados con cinturón y estadísticas. Optimización a 72--90 FPS y háptica menor a 15 ms. Pruebas HEART y corrección de errores. & \textbf{Versión final con los 5 niveles y resultados de las pruebas de usabilidad.}\\ \hline
Entregable 6 & 15 & Presentación & Informe final actualizado, demo en vivo del prototipo y comparación entre el boceto del Hito 1 y lo implementado. Video de respaldo de la demo. & Presentación y demo final del proyecto.\\
}
\end{table}

\textbf{Estado del avance:} a la semana 14 están completados los Entregables 1 a 4 y, del Entregable~5, todo lo que no depende del visor ni de los usuarios: los cinco niveles, la escena de resultados con el cinturón, la preparación del rendimiento y de las pruebas HEART, y las correcciones de la zona segura y de las animaciones. Quedan las mediciones en Meta Quest~2 y las pruebas con usuarios: \pendiente{FPS y latencia medidos en el visor, sesión de 30 minutos y resultados HEART (sección~\ref{sec:e5})}. El detalle de lo implementado está en las secciones~\ref{sec:e2e3}, \ref{sec:e4} y~\ref{sec:e5}.

\subsection{Diagrama de Gantt}
Semanas 10 a 15. Entregables al 15\,\%, 30\,\%, 50\,\%, 70\,\% y 100\,\%. Desarrollo en rojo; pruebas y documentación en dorado; en gris, lo de la semana 14 que espera al visor o a los usuarios. La columna RF / RNF indica los requisitos que se implementan en cada actividad, y la última, su estado al cierre de la semana 14.

\newcommand{\R}{\cellcolor{rojo}}
\newcommand{\G}{\cellcolor{dorado}}
\newcommand{\Pe}{\cellcolor{gray!45}}
\newcommand{\sem}[3]{\parbox[c][1.15cm][c]{0.95cm}{\centering\tiny\textbf{\textcolor{white}{#1}}\\ \textcolor{white}{#2}\\ \textbf{\textcolor{dorado}{#3}}}}
\begin{table}[H]
\centering
\caption{Diagrama de Gantt del proyecto y estado de cada actividad en la semana 14.}
\scriptsize
\setlength{\tabcolsep}{2.5pt}
\renewcommand{\arraystretch}{1.1}
\begin{tabular}{|>{\raggedright\arraybackslash}p{4.3cm}|>{\centering\arraybackslash}p{1.6cm}|*{6}{>{\centering\arraybackslash}p{0.95cm}|}>{\raggedright\arraybackslash}p{2.6cm}|}
\hline
\rowcolor{vino}\hd{Actividad}&\hd{RF / RNF}&\sem{Sem.\ 10}{Entreg.\ 1}{15\,\%}&\sem{Sem.\ 11}{Entreg.\ 2}{30\,\%}&\sem{Sem.\ 12}{Entreg.\ 3}{50\,\%}&\sem{Sem.\ 13}{Entreg.\ 4}{70\,\%}&\sem{Sem.\ 14}{Entreg.\ 5}{100\,\%}&\sem{Sem.\ 15}{Entreg.\ 6}{Present.}&\hd{Estado}\\ \hline
Unity, Meta XR SDK y OVR Camera Rig & RF-01 & \R & & & & & & Completado\\ \hline
Menú principal y paneles de juego y opciones & RF-17, RF-09 & \R & \R & & & & & Completado\\ \hline
Ring y assets del escenario & --- & \R & \R & & & & & Completado\\ \hline
Guantes y colisiones (jab, cross, uppercut) & RF-02, RF-03 & & \R & \R & & & & Completado\\ \hline
Háptica diferenciada & RF-15 & & \R & \R & \R & & & Completado\\ \hline
IA niveles 1 y 2: Glass Joe y Bald Bull & RF-04, RF-05, RF-07, RF-08 & & \R & \R & & & & Completado\\ \hline
Esquive con cabeza y torso & RF-06 & & & \R & & & & Completado\\ \hline
Salud diegética: viñeta y pulso cardíaco & RF-14 & & & \R & \R & & & Completado\\ \hline
Hand Grab Recorder para bloqueos de la IA & --- & & & \R & \R & \Pe & & Hecho; falta grabar la guardia real con el visor\\ \hline
IA niveles 3 y 4: Soda Popinski y Mr.\ Sandman & RF-07, RF-08 & & & & \R & & & Completado\\ \hline
Réferi, cuenta de KO, button mashing y rondas & RF-10 a RF-13 & & & & \R & \R & & Completado (corregido en la semana 14)\\ \hline
IA nivel 5: Iron Mike Tyson & RF-07, RF-08 & & & & & \R & & Hecho; dificultad por confirmar en el visor\\ \hline
Escena de resultados, cinturón y estadísticas & RF-18, RF-19 & & & & & \R & & Completado\\ \hline
Optimización (72--90 FPS, háptica menor a 15 ms) & RNF-01, RNF-02 & & & & \R & \Pe & & Preparada; medición en el visor pendiente\\ \hline
Pruebas de usabilidad HEART y correcciones & RNF-03, RNF-05 & & & & & \Pe & & Protocolo listo; pruebas pendientes\\ \hline
Correcciones tras la prueba en el visor: zona segura, pasos de los personajes, gesto del réferi & RF-12, RF-20, RNF-04 & & & & & \R & & Completado; validar en el visor\\ \hline
Informe final, comparación del boceto y demo & --- & & & & & \G & \G & En curso\\ \hline
Audio espacial 3D (*) & RF-16 & & & \R & \R & & & Completado\\ \hline
Pausa por sistema Guardián (*) & RF-20 & \R & \R & & & \R & & Completado (corregido en la semana 14)\\ \hline
Registro de métricas de sesión (*) & RF-21 & & & & \R & \R & & Completado\\ \hline
\rowcolor{crema}\textbf{AVANCE ACUMULADO} & & \textbf{15\,\%} & \textbf{30\,\%} & \textbf{50\,\%} & \textbf{70\,\%} & \textbf{100\,\%} & \textbf{Present.} & \\ \hline
\end{tabular}
\end{table}
{\small (*) Actividades agregadas al cronograma original para dar cabecera a requisitos funcionales que no tenían una actividad asociada (RF-16, RF-20 y RF-21).}

\subsection{Trazabilidad de requisitos}
La siguiente tabla indica en qué semana y en qué entregable se implementa cada requisito funcional de la sección~3.1 y su estado al cierre de la semana 14. Como el cronograma se organiza por semanas, la semana es la unidad mínima de planificación.

\begin{table}[H]
\centering
\caption{Trazabilidad de requisitos funcionales y estado en la semana 14.}
\label{tab:traza}
\small
\begin{tabularx}{\linewidth}{|c|L|c|c|>{\raggedright\arraybackslash}p{3.1cm}|}
\hline
\cab\hd{ID}&\hd{Requisito (resumen)}&\hd{Semana(s)}&\hd{Entregable}&\hd{Estado}\\ \hline
RF-01 & Vista en primera persona y guantes con seguimiento 6DoF & 10 & 1 & Completado\\ \hline
RF-02 & Detección de colisión guante--oponente (golpe conectado) & 11 & 2 & Completado\\ \hline
RF-03 & Clasificación del golpe: jab, cross o uppercut & 11 & 2 & Completado\\ \hline
RF-04 & Reducción de la salud del oponente según golpe y velocidad & 12 & 3 & Completado\\ \hline
RF-05 & Detección de bloqueo con guardia y reducción de daño & 12 & 3 & Completado\\ \hline
RF-06 & Detección de esquive con cabeza y torso & 12 & 3 & Completado\\ \hline
RF-07 & Cinco oponentes IA con patrones propios & 12 a 14 & 3, 4 y 5 & Completado (5 rivales)\\ \hline
RF-08 & Dificultad ascendente del Nivel 1 al 5 & 12 a 14 & 3, 4 y 5 & Completado en simulación; \pendiente{confirmar la dificultad de Tyson con jugadores}\\ \hline
RF-09 & Desbloqueo del siguiente nivel al derrotar al rival & 11 y 12 & 2 y 3 & Completado (5 niveles)\\ \hline
RF-10 & 3 rounds de 90 s con campana & 13 & 4 & Completado\\ \hline
RF-11 & Ganador por KO o por mayor salud al tercer round & 13 & 4 & Completado\\ \hline
RF-12 & Cuenta del réferi de 10 s y button mashing & 13 y 14 & 4 y 5 & Completado (la cuenta del rival espera a que esté tendido)\\ \hline
RF-13 & Derrota si no supera la cuenta; reintento del nivel & 13 & 4 & Completado\\ \hline
RF-14 & Salud diegética: viñeta periférica y latido, sin HUD & 12 y 13 & 3 y 4 & Completado\\ \hline
RF-15 & Vibración háptica diferenciada & 11 a 13 & 2, 3 y 4 & Completado\\ \hline
RF-16 & Audio espacial 3D (golpes, multitud, réferi) (*) & 12 y 13 & 3 y 4 & Completado\\ \hline
RF-17 & Menú principal diegético: Jugar, Opciones y Salir & 10 y 11 & 1 y 2 & Completado\\ \hline
RF-18 & Resultados por combate: precisión, bloqueos y tiempo & 14 & 5 & Completado\\ \hline
RF-19 & Cinturón de campeón al derrotar al Nivel 5 & 14 & 5 & Completado\\ \hline
RF-20 & Pausa del combate por el sistema Guardián (*) & 10, 11 y 14 & 1, 2 y 5 & Completado; corregido en la semana 14 (sección~\ref{sec:correcciones})\\ \hline
RF-21 & Registro de métricas por sesión (*) & 13 y 14 & 4 y 5 & Completado (incluye la adopción de RNF-03)\\ \hline
\end{tabularx}
\end{table}
{\small (*) Requisitos cuya actividad fue agregada al cronograma (sección 8.2).}

Los requisitos no funcionales no tienen una actividad propia en el cronograma: se verifican con las pruebas de la semana 14. La tabla~\ref{tab:traza-rnf} resume su estado y la evidencia de cada uno (sección~\ref{sec:verif5}).

\begin{table}[H]
\centering
\caption{Estado de los requisitos no funcionales en la semana 14.}
\label{tab:traza-rnf}
\small
\begin{tabularx}{\linewidth}{|c|>{\raggedright\arraybackslash}p{4.2cm}|L|}
\hline
\cab\hd{ID}&\hd{Requisito (resumen)}&\hd{Estado y evidencia}\\ \hline
RNF-01 & 72 FPS; menos del 5\,\% de fotogramas lentos & Medida preparada en el juego (\cs{RegistroSesion}) y protocolo listo. \pendiente{FPS por nivel medidos en Meta Quest 2}\\ \hline
RNF-02 & Háptica en menos de 15 ms & Física a 1/72 s y latencia medida en cada golpe. \pendiente{latencia medida en el visor}\\ \hline
RNF-03 & 80\,\% de novatos con guardia y golpe en 2 min & Tiempos registrados por el juego y protocolo con cronómetro. \pendiente{resultado de las pruebas con usuarios}\\ \hline
RNF-04 & Sin teletransporte ni joystick; desplazamiento físico 1:1 & Completado: no hay teletransporte; la locomoción con joystick de depuración solo funciona en el Editor\\ \hline
RNF-05 & Cinetosis: promedio mínimo de 4 sobre 5 & Cuestionario Likert listo. \pendiente{resultado del cuestionario}\\ \hline
RNF-06 & Área de 1,5 $\times$ 1,5 m y límites del Guardián & Verificado en el Editor: se juega sin moverse del centro y la pausa salta a 5 cm del borde de un área de 1,5 $\times$ 1,5 m. Geometría del Guardián corregida. \pendiente{validar la pausa del Guardián en el visor}\\ \hline
RNF-07 & Al menos el 95\,\% de los impactos registrados & Verificado en el Editor con golpes por física: 68 de 70 (97\,\%) de 1 a 12 m/s\\ \hline
RNF-08 & 30 minutos de sesión continua & Protocolo listo. \pendiente{sesión de 30 minutos en el visor}\\ \hline
RNF-09 & Standalone en Quest 2 y Quest 3 & APK autónomo para Quest 2. Quest 3 no se declara (decisión del equipo); se anota si funciona si hay uno\\ \hline
RNF-10 & Parámetros de cada rival configurables & Completado: un asset \cs{PerfilOponente} por rival\\ \hline
\end{tabularx}
\end{table}

% ============================================================
\section{Referencias Bibliográficas}
\begin{itemize}[label={},leftmargin=0pt,itemsep=5pt]
  \item Slater, M., \& Wilbur, S. (1997). A framework for immersive virtual environments (FIVE): Speculations on the role of presence in virtual environments. \emph{Presence: Teleoperators and Virtual Environments}, 6(6), 603--616.
  \item Bowman, D. A., \& McMahan, R. P. (2007). Virtual reality: How much immersion is enough? \emph{IEEE Computer}, 40(7), 36--43.
  \item Dix, A., Finlay, J., Abowd, G. D., \& Beale, R. (2004). \emph{Human-Computer Interaction} (3rd ed.). Pearson Prentice Hall.
  \item Freytag, G. (1863). \emph{Die Technik des Dramas}. Hirzel.
  \item Meta Platforms. (2024). \emph{Meta XR Core SDK Documentation}. \url{https://developer.oculus.com/documentation/unity/unity-ovr-utilities-overview/}
  \item Unity Technologies. (2022). \emph{Unity 2022 LTS User Manual}. \url{https://docs.unity3d.com/2022.3/Documentation/Manual/}
  \item Google Ventures. (2013). HEART Framework for measuring UX. \url{https://library.gv.com/how-to-choose-the-right-ux-metrics-for-your-product-5f46359ab5be}
  \item Sketchfab. (2024). \emph{3D Model Repository --- Licencias Creative Commons}. \url{https://sketchfab.com}
\end{itemize}

% ============================================================
\section{Implementación del Entregable 3: Avances de la Semana 12 (50\,\% de Avance Acumulado)}
\label{sec:e2e3}
\subsection{Resumen del Hito y Resultado Verificable}
De acuerdo con el cronograma de trabajo (sección~\ref{sec:crono}), el Entregable~3 corresponde a la semana 12 y representa el 50\,\% de avance acumulado. Su resultado verificable exige combates completos y fluidos contra dos rivales (Glass Joe en el Nivel 1 y Bald Bull en el Nivel 2, con la integración técnica adelantada de Soda Popinski en el Nivel 3), donde la salud y el daño se perciban de manera puramente diegética, sin elementos de interfaz gráfica (HUD) sobre la vista del usuario, y donde los oponentes demuestren una inteligencia artificial reactiva con esquemas propios de ataque, bloqueo y desplazamiento.

A continuación se presenta el detalle técnico de las funcionalidades implementadas en Unity 2022 LTS para Meta Quest (OpenXR / Android), organizadas en la arquitectura del proyecto (carpeta \cs{Assets/Scripts}).

\subsection{IA Reactiva de los Rivales: Bloqueo, Velocidad y Acoso (RF-07, RF-08, RNF-10)}
Los contrincantes virtuales dejaron de ser objetivos pasivos y pasaron a constituir agentes autónomos reactivos cuyas decisiones dependen de la conducta del jugador en el cuadrilátero:
\begin{itemize}
  \item \textbf{Bloqueo dinámico por combinaciones (\cs{DefensaOponente}):} cada rival posee una probabilidad base de bloqueo que se incrementa con cada golpe consecutivo que recibe dentro de una ventana de 1,5 s (\cs{ventanaGolpesSeguidos}). De este modo el oponente ``lee'' las combinaciones reiteradas del jugador y sube su guardia con mayor frecuencia hasta un límite calibrado (\cs{bloqueoMaximo}). Tras bloquear, mantiene la guardia durante \cs{segundosGuardia} (0,6 a 0,8 s) protegiéndose de impactos subsecuentes. Al descender los brazos queda expuesto en una ventana de vulnerabilidad (\cs{aberturaTrasGuardia} = 0,6 s) en la que no puede volver a bloquear, premiando la iniciativa táctica del jugador. Mientras ejecuta una animación de ataque, su capacidad de bloqueo se inhabilita por completo, permitiendo contragolpes certeros.
  \item \textbf{Guardia visual mediante cinemática inversa (Hand IK):} en lugar de depender de animaciones grabadas rígidas, la postura defensiva se genera mediante cinemática inversa sobre los avatares humanoides. \cs{DefensaOponente} instancia el componente \cs{RelevoIK} sobre el objeto del \emph{Animator} para capturar el callback nativo \cs{OnAnimatorIK}. Al activarse la guardia, las manos del modelo se posicionan suavemente frente al rostro mediante interpolación cúbica (\cs{Mathf.SmoothStep}), con \emph{IK Pass} activo en la capa base de los controladores de animación. En el Entregable~4 esta posición fija se sustituyó por una trayectoria grabada (sección~\ref{sec:grabadora}).
  \item \textbf{Desplazamiento, acoso y castigo por guardia baja (\cs{AtaqueOponente}):} el rival orienta constantemente su mirada hacia el jugador (\cs{velocidadGiro} = 8 rad/s) y, si el usuario se aleja a una distancia superior a \cs{distanciaCombate} (1,1 a 1,2 m), camina hacia él a \cs{velocidadDesplazamiento} (0,4 a 0,9 m/s), respetando siempre las cuerdas (\cs{LimiteRing.AreaRing}). Si el jugador baja los guantes descuidando su rostro, el temporizador del próximo ataque avanza a ritmo acelerado (\cs{castigoGuardiaBaja}). Tras bloquear un ataque del usuario, programa un contragolpe anticipado (\cs{contraataqueTrasBloqueo}).
  \item \textbf{Freno y colisión de impacto contra la guardia del jugador:} para evitar que los puñetazos del oponente atraviesen visualmente los guantes del usuario al ser bloqueados, se incorporó una lógica de anticipación geométrica (\cs{anticipoGuardia} = 0,08 s). Si el golpe entrante intersecta el volumen cefálico y el jugador mantiene la guardia alta, la animación del rival se detiene (\cs{animator.speed} = 0), permanece frenada contra las manos durante \cs{pausaGuardia} = 0,08 s y se retira con suavidad a la posición neutral.
  \item \textbf{Configuración desacoplada y dificultad ascendente (\cs{PerfilOponente}):} cumpliendo RNF-10, todos los parámetros de combate residen en assets \cs{ScriptableObject} (\cs{Assets/_Project/Data/Oponentes/}), lo que permite calibrar la progresión entre los rivales sin alterar el código. La tabla~\ref{tab:perfiles3} muestra los valores del Entregable~3; la tabla~\ref{tab:perfiles4} los actualiza con los cuatro rivales.
\end{itemize}

\begin{table}[H]
\centering
\caption{Perfiles de los tres primeros rivales en el Entregable 3.}
\label{tab:perfiles3}
\tablita{|L|c|c|c|}{\hd{Parámetro}&\hd{Glass Joe (N1)}&\hd{Bald Bull (N2)}&\hd{Soda Popinski (N3)}}{
Salud máxima (puntos) & 100 & 150 & 200\\ \hline
Probabilidad de bloqueo (base $\to$ máx.) & 8\,\% $\to$ 30\,\% & 20\,\% $\to$ 55\,\% & 30\,\% $\to$ 65\,\%\\ \hline
Multiplicador de velocidad de ataque & $\times$0,90 & $\times$1,15 & $\times$1,25\\ \hline
Velocidad de desplazamiento & 0,4 m/s & 0,7 m/s & 0,9 m/s\\ \hline
Daño de ataque (jab / uppercut) & 4 / 7 & 8 / 14 & 10 / 18\\ \hline
Castigo por guardia baja & +30\,\% & +80\,\% & +100\,\%\\ \hline
Ventana de contragolpe tras bloqueo & 1,0 a 1,6 s & 0,5 a 1,0 s & 0,4 a 0,8 s\\
}
\end{table}

\subsection{Sistema Defensivo del Jugador: Bloqueo y Esquive Volumétrico (RF-05, RF-06)}
El componente \cs{SaludJugador}, situado en el \cs{OVRCameraRig}, evalúa cada golpe entrante del oponente mediante geometría espacial en tiempo real, resolviendo entre tres desenlaces posibles:
\begin{itemize}
  \item \textbf{Esquive físico con cabeza y torso (RF-06):} el puñetazo del rival traza una trayectoria rectilínea desde el origen de su puño hasta el alcance configurado (\cs{alcanceAtaque}). El script calcula el punto más cercano de ese segmento al centro de visión del visor (\cs{centerEyeAnchor}). Si el usuario desplaza su cabeza o ladea el torso a una distancia radial mayor a \cs{radioImpacto} = 0,3 m de la línea de ataque, el golpe se resuelve como esquivado, sin pérdida de salud.
  \item \textbf{Bloqueo con guantes (RF-05):} si el ataque intersecta el volumen cefálico, el sistema comprueba la postura de los brazos: hay guardia activa si ambos controladores están a una distancia máxima de 0,4 m de la cabeza (\cs{distanciaGuardia}) y en el hemisferio frontal respecto a la orientación de los ojos. Al bloquear con éxito, el jugador absorbe únicamente el 25\,\% del daño (\cs{danoTrasBloqueo} = 0,25), con un sonido de choque de cuero y una respuesta háptica en ambos mandos.
  \item \textbf{Impacto pleno:} si no esquiva ni bloquea, la salud se reduce en la totalidad del daño asignado al ataque, con la vibración de impacto directo y el sonido 3D correspondiente.
\end{itemize}

\subsection{Salud Diegética y Retroalimentación Sensorial sin HUD (RF-14, RF-15)}
En cumplimiento de RF-14, la experiencia prescinde de barras de vida, indicadores numéricos o elementos de pantalla flotantes. La condición física se transmite a través de los canales sensoriales del usuario:
\begin{itemize}
  \item \textbf{Viñeta perimetral procedural (\cs{FeedbackSalud}):} para evitar el coste del posprocesado en el hardware móvil de Meta Quest, el componente genera en tiempo de ejecución una textura de 128 $\times$ 128 píxeles con un perfil cuadrático suave, transparente en el centro y opaco en los bordes. Montada sobre un Canvas en \emph{World Space} a 0,3 m de los ojos, su opacidad responde a la salud normalizada: a mayor daño, más se estrecha el campo visual periférico. Con menos de 20 puntos de salud se activa un oscurecimiento intermitente suave (\cs{periodoParpadeo} = 1,6 s) que emula la pérdida inminente de consciencia.
  \item \textbf{Pulso cardíaco adaptativo:} se sintetizó un clip de un solo latido (lub-dub). El ritmo de reproducción se modula entre 60 y 90 BPM según la severidad del daño. El volumen y el ritmo tienen curvas asimétricas: se aceleran de inmediato al recibir castigo y desaceleran pausadamente al reponerse, sin cortes abruptos.
  \item \textbf{Háptica diferenciada ampliada (RF-15):} se incorporaron perfiles únicos para cada evento táctil (tabla~\ref{tab:haptica3}).
\end{itemize}

\begin{table}[H]
\centering
\caption{Patrones de vibración del Entregable 3.}
\label{tab:haptica3}
\tablita{|>{\raggedright\arraybackslash}p{4.2cm}|L|L|}{\hd{Evento de combate}&\hd{Patrón en los mandos Touch}&\hd{Sensación}}{
Jab conectado (jugador) & 1 pulso de 0,7 durante 50 ms. & Seco, rápido y penetrante.\\ \hline
Cross conectado (jugador) & 1 pulso de 1,0 durante 140 ms. & Contundente, pesado y prolongado.\\ \hline
Uppercut conectado (jugador) & Pulso 0,5 (50 ms), pausa 40 ms y pulso 1,0 (100 ms). & Doble impulso ascendente.\\ \hline
Gancho lateral conectado & 1 pulso de 0,35 durante 60 ms. & Moderado y rasante.\\ \hline
Golpe bloqueado por el rival & Doble pulso de baja frecuencia (25\,\%, amplitudes 0,8 y 0,5). & Toque sordo y apagado de absorción.\\ \hline
Golpe recibido por el jugador & Pulso de 1,0 durante 250 ms en ambos controles. & Sacudida de impacto directo a la cabeza.\\ \hline
Ataque rival bloqueado por el jugador & Doble pulso de 0,6 durante 50 ms con pausa de 40 ms. & Resistencia sólida en ambos antebrazos.\\
}
\end{table}

\subsection{Paisaje Sonoro Espacial 3D (RF-16)}
\begin{itemize}
  \item \textbf{Síntesis de audio procedural:} el script \cs{Tools/sintetizar_audio.py} genera los efectos en WAV mono de 16 bits (22\,050 Hz) mediante filtros Biquad: golpe seco, chasquido de guantes, impacto bloqueado, latido y murmullo de multitud.
  \item \textbf{Ambiente de multitud envolvente (\cs{AmbienteMultitud}):} cada escena de combate tiene un objeto central que instancia cuatro fuentes de audio 3D en anillo a 6 m del centro y 2,5 m de altura, con el murmullo de las gradas en bucle, atenuación logarítmica y desfases independientes. Cuando el oponente o el jugador cae, se reproduce una ovación desde las dos gradas más cercanas al incidente.
  \item \textbf{Audio posicional de impactos y bloqueos:} cada golpe emite su sonido desde el punto de contacto exacto (\cs{Collider.ClosestPoint}), de modo que el usuario ubica si el impacto fue en la sien, la mandíbula o el plexo solar.
\end{itemize}

\subsection{Estructura de Progresión, Niveles y Mecánicas Adelantadas (RF-09, RF-12, RF-13)}
\begin{itemize}
  \item \textbf{Integración de los niveles 1, 2 y 3:} \cs{Scene_Ring_Level1} (Glass Joe), \cs{Scene_Ring_Level2} (Bald Bull, escala 1,15) y \cs{Scene_Ring_Level3} (Soda Popinski, con esqueleto de Mixamo y clips de guardia). Los dos personajes adicionales usan un \cs{AnimatorOverrideController} sobre la máquina de estados base de Glass Joe, reutilizando transiciones y lógica de caídas.
  \item \textbf{Persistencia y desbloqueo secuencial (RF-09):} al vencer a un rival, \cs{ControlCombate} guarda el nivel alcanzado en \cs{PlayerPrefs} (clave \cs{"nivelSuperado"}). El menú de selección inhabilita los botones bloqueados, muestra el rótulo BLOQUEADO y los atenúa.
  \item \textbf{Cartel de victoria y transición:} al producirse la derrota definitiva del rival, \cs{ControlCombate} espera 1,5 s para dejar ver el derribo, despliega un panel en \emph{World Space} (``¡VICTORIA!'') frente a la cámara y carga el siguiente combate o regresa al menú principal.
  \item \textbf{Cuenta de 10 s y recuperación por button mashing (adelantado del Entregable~4):} cuando la salud del jugador llega a cero, \cs{RecuperacionJugador} inicia una cuenta de 10 s en la que pulsar A, B, X o Y lo levanta. Las oportunidades son limitadas y decrecientes (\cs{saludPorCaida} = \{40, 30, 20\}); una cuarta caída produce el K.O. directo y la derrota. En el Entregable~4 este mecanismo se completó con el réferi (sección~\ref{sec:referi}).
\end{itemize}

\subsection{Seguridad en Realidad Virtual y Confinamiento en el Ring (RF-20)}
\begin{itemize}
  \item \textbf{Pausa automática por el Guardián (\cs{PausaGuardian}):} monitorea en cada fotograma la geometría del área jugable del visor (\cs{OVRBoundary.GetGeometry}). Si la cabeza cruza el límite seguro (con un margen de 0,1 m), o si la aplicación pierde el foco de entrada por apertura del menú universal o retiro del visor, el combate se pausa de inmediato mostrando ``VUELVE AL ÁREA SEGURA'' y se reanuda al regresar.
  \item \textbf{Confinamiento en el cuadrilátero (\cs{LimiteRing}):} impide que el jugador traspase las cuerdas restringiendo la posición horizontal del \cs{OVRCameraRig}. Incorpora además un radio de repulsión respecto al rival (\cs{distanciaMinima}), de modo que no pueda interpenetrar su torso ni su cabeza.
\end{itemize}

\subsection{Resumen de Cumplimiento (Semana 12)}
\begin{table}[H]
\centering
\caption{Cumplimiento del Entregable 3.}
\tablita{|>{\raggedright\arraybackslash}p{4.2cm}|L|c|}{\hd{Actividad planificada}&\hd{Implementación en Unity / scripts C\#}&\hd{Estado}}{
IA reactiva Nivel 1 (Glass Joe) y Nivel 2 (Bald Bull) & \cs{AtaqueOponente}, \cs{DefensaOponente}, \cs{RelevoIK} y perfiles \cs{ScriptableObject}: desplazamiento, acoso, bloqueo dinámico y contragolpe. & \ok\\ \hline
Esquive con cabeza y torso (RF-06) & \cs{SaludJugador}: distancia a la línea del golpe; esquive sin daño a más de 0,3 m. & \ok\\ \hline
Bloqueo del jugador con reducción de daño (RF-05) & \cs{SaludJugador}: guantes frente a la cara a 0,4 m o menos; absorbe el 75\,\% del daño. & \ok\\ \hline
Salud diegética con viñeta y pulso cardíaco (RF-14) & \cs{FeedbackSalud}: viñeta procedural, parpadeo crítico y latido de 60 a 90 BPM. Sin HUD. & \ok\\ \hline
Háptica diferenciada (RF-15) & \cs{Puno} y \cs{SaludJugador}: patrones para jab, cross, uppercut, gancho, bloqueo del rival y daño recibido. & \ok\\ \hline
Audio espacial 3D de golpes y multitud (RF-16) & \cs{AmbienteMultitud} con 4 gradas y ovación por cercanía a las caídas; efectos WAV sintetizados. & \ok\\ \hline
Desbloqueo de niveles y progresión (RF-09) & \cs{ControlCombate} y \cs{MainMenu}: \cs{PlayerPrefs}, botones bloqueados y cartel de victoria. & \ok\\ \hline
Pausa por el Guardián (RF-20) & \cs{PausaGuardian}: pausa al salir del área de juego o perder el foco, con aviso. & \ok\\ \hline
Cuenta de 10 s y button mashing (RF-12, RF-13) & \cs{RecuperacionJugador}: cuenta y 3 oportunidades con salud decreciente. & \ok (adelanto)\\ \hline
Resultado verificable del hito (50\,\%) & Combates completos y fluidos contra dos rivales (y el Nivel 3 jugable), con salud diegética sin HUD. & Cumplido\\
}
\end{table}

\subsection{Referencias Bibliográficas Adicionales}
\begin{itemize}[label={},leftmargin=0pt,itemsep=5pt]
  \item Begault, D. R. (1994). \emph{3-D Sound for Virtual Reality and Multimedia}. Academic Press Professional.
  \item Marr, D. (1982). \emph{Vision: A computational investigation into the human representation and processing of visual information}. W. H. Freeman.
  \item Meta Platforms. (2024). \emph{OVRBoundary and Spatial Audio in Meta XR SDK}. Oculus Developer Documentation.
  \item Unity Technologies. (2022). \emph{Inverse Kinematics in Humanoid Avatars}. Unity 2022 LTS User Manual.
  \item ISO/IEC 23005-6:2021. \emph{Information technology --- Media context and control --- Part 6: Common types and tools}.
\end{itemize}

% ============================================================
\section{Implementación del Entregable 4: Avances de la Semana 13 (70\,\% de Avance Acumulado)}
\label{sec:e4}
\subsection{Resumen del Hito y Resultado Verificable}
Según el cronograma (sección~\ref{sec:crono}), el Entregable~4 comprende la IA de los niveles 3 y 4 (Soda Popinski y Mr.\ Sandman), el réferi con cuenta regresiva y \emph{button mashing} tras el KO, las rondas de 3 $\times$ 90 s, el Hand Grab Recorder en los bloqueos y la háptica diferenciada. Su resultado verificable es: \textbf{cuatro rivales jugables y ciclo de KO funcionando}. Se cumplió: el jugador recorre los niveles 1 a 4 en una partida continua, cada combate se organiza en rounds con campana, cualquiera de los dos contendientes puede caer y ser contado por el réferi, y el combate termina por KO o por salud al final del tercer round.

Además se adelantó el registro de métricas por sesión (RF-21, previsto para las semanas 13 y 14), que alimentará las pruebas HEART de la semana 14.

Todo se implementó en Unity 2022 LTS para Meta Quest (OpenXR, Android). El réferi y Mr.\ Sandman usan sus modelos 3D definitivos, riggeados por el equipo; Mr.\ Sandman reutiliza las animaciones de Soda Popinski gracias al avatar humanoide de Unity. El modelo de Iron Mike Tyson también está listo y se integrará con el nivel 5.

\begin{table}[H]
\centering
\caption{Bloques de trabajo del Entregable 4.}
\tablita{|c|L|>{\raggedright\arraybackslash}p{2.6cm}|c|}{\hd{Bloque}&\hd{Contenido}&\hd{Requisitos}&\hd{Sección}}{
A & Rondas de 3 $\times$ 90 s con campana y ganador por KO o por salud & RF-10, RF-11 & \ref{sec:rondas}\\ \hline
B y F & Réferi: cuenta de 10 s en voz 3D, button mashing sin HUD y réferi visible & RF-12, RF-13, RF-16 & \ref{sec:referi}\\ \hline
C & IA de los niveles 3 y 4: combos de Soda, esquive y contraataque de Sandman & RF-07, RF-08, RF-09 & \ref{sec:ia}\\ \hline
D & Hand Grab Recorder: guardias grabadas con el visor para los bloqueos & RF-07 (presentación) & \ref{sec:grabadora}\\ \hline
E & Háptica diferenciada completa & RF-15 & \ref{sec:haptica}\\ \hline
G & Registro de métricas de la sesión (JSON y CSV) & RF-21 & \ref{sec:registro}\\
}
\end{table}

\subsection{Rondas de 3 $\times$ 90 Segundos (RF-10, RF-11)}
\label{sec:rondas}
El script \cs{ControlRondas}, en un objeto propio de cada escena de pelea, organiza el combate en las fases preparación, round, descanso y final. Una campana espacial 3D en una esquina del ring suena una vez al empezar cada round y tres veces al terminarlo, para distinguir ambos momentos solo de oído.

\begin{table}[H]
\centering
\caption{Parámetros de las rondas.}
\tablita{|L|c|>{\ttfamily\small}l|}{\hd{Parámetro}&\hd{Valor}&\hd{Campo}}{
Rounds & 3 & rounds\\ \hline
Duración del round & 90 s & segundosPorRound\\ \hline
Descanso entre rounds & 20 s & segundosDescanso\\ \hline
Preparación antes del primer round & 5 s & segundosPreparacion\\ \hline
Salud que recupera el jugador en cada descanso & 20\,\% de la máxima & fraccionSaludDescanso\\
}
\end{table}

\begin{itemize}
  \item \textbf{Cronómetro:} usa tiempo escalado, así que se detiene con el menú de pausa y con la pausa del Guardián, y no avanza mientras el jugador está en la lona. Un round que llega a cero espera mientras haya una cuenta del réferi en marcha.
  \item \textbf{Fuera de los rounds no se pelea:} el rival no ataca y los golpes del jugador no cuentan, para que no se le pueda castigar gratis durante el descanso. En cada descanso el rival se levanta, vuelve a su esquina y baja la guardia.
  \item \textbf{Resultado:} un KO termina el combate al instante. Si se completan los tres rounds, gana quien conserve más salud en porcentaje. \textbf{En caso de empate gana el jugador}, para no alargar un combate que el diseño limita a tres rounds por la fatiga física (sección~7.2).
  \item \textbf{Duración del descanso:} el diseño no la fijaba. Se eligieron 20 s para no cortar el ritmo de una sesión en RV; se cambia en el componente sin tocar el código.
\end{itemize}

Al ganar, \cs{ControlCombate} guarda el nivel superado y carga la escena siguiente; al perder, la partida vuelve al menú con el nivel todavía abierto para reintentarlo (RF-13).

\subsection{Réferi, Cuenta Regresiva y Button Mashing (RF-12, RF-13, RF-16)}
\label{sec:referi}
El réferi (script \cs{Arbitro}) cuenta en voz alta del 1 al 10 con una voz espacial 3D que sale de su cabeza. Los diez números se sintetizaron con un sintetizador de formantes (\cs{Tools/sintetizar_audio.py}) y se pueden sustituir por grabaciones reales sin tocar el código.

\begin{itemize}
  \item \textbf{Caída del jugador:} al quedarse sin salud empieza una cuenta de 10 s, y pulsar rápido los botones A, B, X o Y lo levanta. Cada pulsación aclara la visión periférica, calma el latido y vibra más fuerte (de 0,25 a 1,0); al levantarse hay dos vibraciones fuertes. Si la cuenta termina, el jugador pierde por K.O. y vuelve al menú con el nivel disponible (RF-13). No hay panel con números ni barras: se mantiene el principio de no usar HUD.
  \item \textbf{Caída del rival:} el réferi cuenta y el rival se levanta al terminar la cuenta (\cs{cuentaKnockdown} = 10 en los tres rivales de ese momento). La animación de caída se sostiene en el piso hasta entonces. Llegar a 10 con el rival no es K.O.: su derrota la deciden las tres caídas por umbrales de salud, y la \textbf{caída final no se cuenta} porque termina el combate. Mientras el rival se incorpora, durante 1,3 s, los golpes no cuentan.
  \item \textbf{Ritmo y pausa:} cada número suena al empezar su segundo, así que el 10 suena en $t = 9$ s y la cuenta acaba en $t = 10$ s. Como un \cs{AudioSource} ignora \cs{Time.timeScale}, el réferi pausa y reanuda la voz a mano al pausar el juego.
\end{itemize}

\textbf{Réferi visible.} El script \cs{CuerpoArbitro} da cuerpo al réferi (figuras~\ref{fig:referi} y~\ref{fig:cuentas}). Se mueve por el ring, gira y marca la cuenta con el brazo, como en el boceto del estado de knockout (figura~\ref{fig:boceto-ko}). Se desliza, porque el modelo no trae animación de caminar. Nunca se acerca a menos de 0,6 m de un boxeador y se queda dentro de las cuerdas.

\begin{table}[H]
\centering
\caption{Posición y orientación del réferi según la fase.}
\tablita{|>{\raggedright\arraybackslash}p{3.2cm}|L|>{\raggedright\arraybackslash}p{3.6cm}|}{\hd{Momento}&\hd{Dónde se pone}&\hd{Hacia dónde mira}}{
Round & A 2 m de la línea entre el jugador y el rival, medio agachado en postura de vigilar, en la periferia de la vista & Al centro de la pelea\\ \hline
Cuenta del rival & Junto al rival, del lado del jugador (0,7 m hacia el jugador y 0,7 m a un lado) & Al rival, al que señala con el brazo\\ \hline
Cuenta del jugador & Delante del jugador (1,05 m) y a un lado (0,4 m) para no superponerse al rival, del lado que lo deja más dentro del ring & Al jugador\\ \hline
Preparación, descanso y final & Su sitio inicial, en un costado & Al centro del ring\\
}
\end{table}

En la cuenta tiene el brazo arriba; con cada número lo baja de golpe en 0,12 s y lo vuelve a subir en 0,45 s, sincronizado con la voz. Sobre su cabeza, y solo en la cuenta del jugador, aparece el cartel ``¡ARRIBA! PULSA A B X Y'', que indica cómo levantarse. Está anclado al réferi dentro del mundo y no a la cámara, de modo que no es un HUD.

\begin{figure}[H]
  \centering
  \begin{subfigure}[t]{0.48\linewidth}
    \includegraphics[width=\linewidth]{02_nivel1_glass_joe.png}
    \caption{Nivel 1: el réferi vigila a la izquierda de la pelea.}
  \end{subfigure}\hfill
  \begin{subfigure}[t]{0.48\linewidth}
    \includegraphics[width=\linewidth]{06_referi.png}
    \caption{Primer plano del réferi en postura de vigilar.}
  \end{subfigure}
  \caption{El réferi durante el round.}
  \label{fig:referi}
\end{figure}

\begin{figure}[H]
  \centering
  \begin{subfigure}[t]{0.48\linewidth}
    \includegraphics[width=\linewidth]{07_cuenta_del_jugador.png}
    \caption{Cuenta del jugador: réferi delante, brazo en alto y cartel de los botones.}
  \end{subfigure}\hfill
  \begin{subfigure}[t]{0.48\linewidth}
    \includegraphics[width=\linewidth]{08_cuenta_del_rival.png}
    \caption{Cuenta del rival: el réferi marca la cuenta junto a Mr.\ Sandman en la lona.}
  \end{subfigure}
  \caption{Las dos cuentas del réferi, desde los ojos del jugador.}
  \label{fig:cuentas}
\end{figure}

\textbf{Modelo del réferi.} El modelo 3D (camisa azul, pantalón oscuro y pajarita) tiene un esqueleto humanoide y dos poses: de pie en T y agachado vigilando. Su malla original tenía 378\,936 triángulos, unas 50 veces más que la de los rivales (6\,652 la de Glass Joe y 8\,636 la de Bald Bull), demasiado para el Meta Quest 2. En un \emph{SkinnedMeshRenderer} el coste es por vértice en cada fotograma, así que el \emph{Mesh LOD} de Unity 2022 no servía: comparte el buffer de vértices. El script \cs{Tools/reducir_referi.py} aplica en Blender una decimación antes del modificador de armadura, que conserva los pesos de los huesos, el UV y el material, y deja la malla en 20\,000 triángulos. A partir de la pose en T se generó una pose de pie relajada que respira, y un \emph{blend tree} lleva al réferi suavemente de una postura a otra según la fase del combate. El gesto de la cuenta se aplica sobre el brazo derecho del esqueleto después de la animación, así que sirve para cualquier modelo humanoide. La herramienta \emph{Tools > Boxeo > Montar réferi} lo monta en las cuatro escenas de combate.

\subsection{IA de los Niveles 3 y 4 (RF-07, RF-08)}
\label{sec:ia}
Soda Popinski lanza combinaciones de ganchos y Mr.\ Sandman esquiva y contraataca. Los valores de cada rival viven en su asset \cs{PerfilOponente} (RNF-10).

\subsubsection{Combos de ataque}
\cs{AtaqueOponente} lanza en cada ataque la secuencia \cs{combo} del perfil; con la secuencia vacía sigue el ciclo anterior (jabs y un uppercut periódico), así que Glass Joe y Bald Bull no cambian.
\begin{itemize}
  \item \textbf{Encadenado:} el siguiente golpe arranca \cs{segundosEntreGolpesCombo} (0,12 s) después del impacto del anterior, sin volver a la guardia. Cada golpe vuelve a apuntar a la cabeza del jugador, así que moverse esquiva un golpe, no todo el combo.
  \item \textbf{Qué lo corta:} un golpe conectado al rival, una caída, el jugador en el piso, la pausa del round o desactivar el componente. Cortar borra el resto del combo y el impacto pendiente, sin dejar impactos fantasma.
  \item \textbf{Qué no lo corta:} un golpe que el jugador bloquea, porque el rival sigue presionando sobre la guardia.
  \item \textbf{Entre dos golpes} el rival no bloquea ni esquiva, de modo que contragolpear a mitad del combo siempre entra.
\end{itemize}

\textbf{Soda Popinski (nivel 3).} Su combo es izquierda, derecha, izquierda, con la velocidad de ataque en $\times$1,3. Los clips de gancho hacían avanzar al personaje 1,33 y 0,71 m, lo que lo metería en el jugador; se quitó el desplazamiento en la importación y el cuerpo se queda en su sitio. Se midió el momento de impacto de cada clip a 30 fotogramas por segundo (tabla~\ref{tab:impactos}).

\begin{table}[H]
\centering
\caption{Tiempos de impacto medidos en los clips de ataque.}
\label{tab:impactos}
\tablita{|L|c|c|c|c|}{\hd{Clip (a escala 1)}&\hd{Impacto en el clip}&\hd{Altura del puño}&\hd{Soda $\times$1,3}&\hd{Sandman $\times$1,4}}{
Gancho izquierdo & 0,50 s & 1,93 m (cabeza 2,04 m) & 0,38 s & 0,36 s\\ \hline
Gancho derecho & 0,80 s & 1,62 m (cabeza 1,84 m) & 0,62 s & no lo usa\\ \hline
Barrido de \cs{Ataque_Fuerte}, primer golpe & 0,30 s & 1,34 m (cabeza 1,82 m) & no lo usa & 0,21 s\\
}
\end{table}

\textbf{Mr.\ Sandman (nivel 4).} Su combo es gancho, barrido, gancho, con velocidad de ataque $\times$1,4 y la espera entre ataques más corta (2,8 a 4 s). Además esquiva con el componente \cs{EsquiveOponente}:
\begin{itemize}
  \item \textbf{Decisión al contacto,} antes que el bloqueo: solo al entrar el guante en una zona de golpe se sabe que hay un golpe de verdad, con su tipo y su zona, así que un amago no lo hace saltar. La probabilidad es del 45\,\% con 1,0 s de enfriamiento.
  \item \textbf{Orden con el bloqueo:} con la guardia arriba, bloquea y no esquiva; en la abertura tras bajar la guardia no hace ninguna de las dos cosas, y es la ventana del jugador; si no, tira el esquive y, si falla, tira el bloqueo.
  \item \textbf{Movimiento:} se aparta 0,35 m hacia el lado contrario del guante o, en el 30\,\% de los casos, hacia atrás, sin salir del ring, y vuelve a su sitio. No esquiva mientras ataca, cae o se incorpora. Un golpe esquivado no quita salud, no suena, no vibra con fuerza (ver sección~\ref{sec:haptica}) y no cuenta como golpe seguido para \cs{DefensaOponente}.
  \item \textbf{Contraataque:} tras esquivar, su siguiente ataque llega entre 0,25 y 0,5 s después, igual que tras un bloqueo.
\end{itemize}

\textbf{Modelo definitivo de Mr.\ Sandman.} El equipo riggeó el modelo de Sketchfab en Blender: esqueleto humanoide de 21 huesos con nombres propios, en pose T y con 1,88 m. No es un rig de Mixamo y no trae animaciones. \cs{Tools/exportar_sandman.py} abre el archivo de Blender sin interfaz, une las 22 piezas en una sola malla (en el Quest cada \emph{SkinnedMeshRenderer} se calcula aparte), guarda las texturas y exporta el FBX de 14\,131 triángulos, entre los rivales y el réferi. En Unity usa las animaciones de Soda retargeteadas por el avatar humanoide, así que los tiempos de impacto medidos siguen valiendo.

Como el retarget humanoide coloca el cuerpo por su centro de masa y no sabe cuánto calzado hay bajo el tobillo, con las animaciones de Soda las botas de Sandman (22 cm bajo el tobillo) quedaban hundidas 23 cm en la lona. El componente \cs{ApoyoEnSuelo} corre después de la animación, mide la posición de las suelas (calculada en el Editor y guardada, para que la malla no necesite lectura en el Quest) y sube la cadera lo que se hayan hundido. Nunca la baja, por lo que no afecta al rival tendido en la lona. La figura~\ref{fig:sandman} muestra el resultado.

\begin{figure}[H]
  \centering
  \includegraphics[width=0.7\linewidth]{05_nivel4_mr_sandman.png}
  \caption{Mr.\ Sandman con su modelo definitivo, con las botas sobre la lona.}
  \label{fig:sandman}
\end{figure}

\subsubsection{Parámetros de los cuatro rivales y dificultad ascendente}
\begin{table}[H]
\centering
\caption{Perfiles de los cuatro rivales (actualiza la tabla~\ref{tab:perfiles3}).}
\label{tab:perfiles4}
\small
\begin{tabularx}{\linewidth}{|>{\raggedright\arraybackslash}p{3.1cm}|C|C|C|C|}
\hline
\cab\hd{Parámetro}&\hd{Glass Joe}&\hd{Bald Bull}&\hd{Soda Popinski}&\hd{Mr.\ Sandman}\\ \hline
Salud & 100 & 150 & 200 & 175\\ \hline
Bloqueo (base $\to$ máx.) & 8\,\% $\to$ 30\,\% & 20\,\% $\to$ 55\,\% & 30\,\% $\to$ 65\,\% & 25\,\% $\to$ 55\,\%\\ \hline
Esquive & no & no & no & 45\,\%, 1,0 s de enfriamiento, 0,35 m\\ \hline
Velocidad de ataque & $\times$0,9 & $\times$1,15 & $\times$1,3 & $\times$1,4\\ \hline
Desplazamiento & 0,4 m/s & 0,7 m/s & 0,9 m/s & 1,1 m/s\\ \hline
Patrón & jab $\times$3, uppercut & jab $\times$2, uppercut & izquierda, derecha, izquierda & gancho, barrido, gancho\\ \hline
Pausa entre golpes del combo & --- & --- & 0,12 s & 0,12 s\\ \hline
Espera entre ataques & 4 a 6 s & 3,5 a 5 s & 3,2 a 4,6 s & 2,8 a 4 s\\ \hline
Daño por golpe & 4 / 7 & 8 / 14 & 9 & 8 / 12\\ \hline
Castigo por guardia baja & +30\,\% & +80\,\% & +100\,\% & +100\,\%\\ \hline
\end{tabularx}
\end{table}

Mr.\ Sandman tiene menos salud que Soda, como pide la salud intermedia del diseño. Lo compensan el esquive y un ritmo más alto, y cada uno de sus golpes quita menos. Para comprobar que la dificultad sube del nivel 1 al 4 (RF-08) se usaron dos medidas (tabla~\ref{tab:dificultad}): un modelo de 20\,000 peleas simuladas que copia la lógica de golpe, bloqueo y esquive, y una pelea en Play por rival con un jugador programado (ráfagas de 5 golpes y guardia siempre arriba, a velocidad $\times$2). Son estimaciones para comparar rivales, no una medida del balance real.

\begin{table}[H]
\centering
\caption{Dificultad relativa de los cuatro rivales.}
\label{tab:dificultad}
\small
\begin{tabularx}{\linewidth}{|>{\raggedright\arraybackslash}p{4.4cm}|C|C|C|C|}
\hline
\cab\hd{Medida}&\hd{Glass Joe}&\hd{Bald Bull}&\hd{Soda Popinski}&\hd{Mr.\ Sandman}\\ \hline
Daño por segundo si todo conecta, guardia arriba & 0,95 & 2,35 & 5,24 & 6,65\\ \hline
Golpes para ganarle, 1 golpe cada 0,5 s (modelo) & 19 & 36 & 51 & 52\\ \hline
Golpes para ganarle, 1 golpe cada 1,5 s (modelo) & 18 & 31 & 46 & 62\\ \hline
Pelea simulada: segundos de juego & 30 & 64 & 83 & 95\\ \hline
Pelea simulada: salud final del jugador & 99 & 85 & 39 & 7\\ \hline
\end{tabularx}
\end{table}

\subsection{Hand Grab Recorder en los Bloqueos}
\label{sec:grabadora}
Para que los bloqueos de los rivales no sean un movimiento rígido en línea recta, se implementó un grabador de guardias:
\begin{itemize}
  \item \textbf{Grabación (\cs{GrabadoraGuardia}):} con el visor puesto y las manos abajo, el jugador pulsa a la vez los dos sticks. Tras una cuenta atrás háptica de tres tics, el script registra durante 2 s, a 60 muestras por segundo, la posición y la rotación de ambas manos respecto a la cabeza, y conserva solo el tramo en que se sube la guardia. El resultado se guarda como JSON y como asset \cs{GrabacionGuardia}. La herramienta \emph{Tools > Boxeo > Importar guardias grabadas} convierte en assets los JSON sacados del Quest.
  \item \textbf{Reproducción:} cuando un rival bloquea, \cs{DefensaOponente} lleva sus guantes por esa trayectoria mediante cinemática inversa de manos, hacia delante al subir la guardia y al revés al bajarla. La guardia usada por cada rival es el campo \cs{guardiaGrabada} de su perfil (RNF-10).
\end{itemize}

Decisiones de diseño:
\begin{itemize}
  \item \textbf{Marco de la cabeza sin inclinación:} las manos se guardan respecto a la cabeza con la vertical del mundo, que es el mismo marco con el que el rival, que siempre mira al jugador, coloca los guantes. Inclinar la cabeza al grabar no tuerce la guardia.
  \item \textbf{Solo la subida:} el recorte empieza cuando una mano supera 0,25 m/s y termina cuando, pasado el pico, las dos bajan de 0,15 m/s durante 4 muestras.
  \item \textbf{Parte de la pose animada:} el jugador graba con las manos en el pecho, pero los rivales ya tienen las manos arriba en reposo, así que reproducirla tal cual hundía los guantes bajo la pose animada. La trayectoria arranca en la pose animada y esa diferencia se borra a lo largo de la subida: se conserva la forma del movimiento y no hay saltos.
  \item \textbf{Escala por altura:} la trayectoria se escala con la altura de la cabeza del rival sobre la de quien grabó, no con la escala del objeto. Con la guardia fija anterior, los guantes de Soda Popinski (modelo grande a escala 0,8) atravesaban su cara; la trayectoria escalada lo corrige (figura~\ref{fig:guardias}).
  \item \textbf{Velocidad:} se mantienen 0,06 s para subir y 0,3 s para bajar. Un bloqueo se decide al contacto y la guardia tiene que estar arriba al instante, así que la subida grabada se comprime; el movimiento grabado se aprecia sobre todo al bajar.
\end{itemize}

\begin{figure}[H]
  \centering
  \begin{subfigure}[t]{0.40\linewidth}
    \includegraphics[width=\linewidth]{09_guardia_fija.png}
    \caption{Guardia fija anterior.}
  \end{subfigure}\hfill
  \begin{subfigure}[t]{0.40\linewidth}
    \includegraphics[width=\linewidth]{10_guardia_grabada.png}
    \caption{Trayectoria grabada y escalada.}
  \end{subfigure}
  \caption{Guardia de Soda Popinski: los guantes ya no atraviesan la cara.}
  \label{fig:guardias}
\end{figure}

\subsection{Háptica Diferenciada Completa (RF-15)}
\label{sec:haptica}
Hasta el Entregable~3 había un patrón por tipo de golpe conectado, el golpe bloqueado por el rival, el golpe recibido y el ataque bloqueado por el jugador (tabla~\ref{tab:haptica3}). En este entregable se completó el repertorio para que el jugador distinga cada evento sin mirar (tabla~\ref{tab:haptica4}).

\begin{table}[H]
\centering
\caption{Patrones hápticos añadidos en el Entregable 4.}
\label{tab:haptica4}
\tablita{|>{\raggedright\arraybackslash}p{4cm}|L|L|}{\hd{Evento}&\hd{Patrón}&\hd{Sensación}}{
El rival esquiva el golpe & 1 pulso de 0,15 durante 30 ms & Roce leve: el puño no tocó nada sólido\\ \hline
El golpe derriba al rival & 3 pulsos de 1,0 (80, 80 y 220 ms), con 40 ms entre ellos & Recompensa clara, distinta de cualquier golpe\\ \hline
Cada número de la cuenta del jugador & Pulso de 30 ms en ambos controles, de 0,2 (con el 1) a 0,6 (con el 10) & Marca el paso del tiempo e invita a pulsar\\ \hline
Pulsación de A, B, X o Y en la cuenta & Pulso de 50 ms, de 0,25 a 1,0 & El avance hacia levantarse\\ \hline
Levantarse & 2 pulsos fuertes & Ya no hace falta pulsar\\
}
\end{table}

\begin{itemize}
  \item \textbf{Esquivado:} sin ninguna señal, el jugador no distinguiría un esquive de un fallo del seguimiento. El roce es casi imperceptible y no se confunde con ningún impacto.
  \item \textbf{Derribo:} es el momento de más recompensa del combate. Su patrón es más largo que cualquier golpe y se reconoce sin mirar.
  \item \textbf{Tic de la cuenta:} sin panel, nada invitaba a pulsar. El tic va al ritmo de la voz, se nota aunque el jugador no mire al réferi y sube de fuerza con la urgencia.
  \item \textbf{El esquive del jugador no vibra:} apartar la cabeza ya lo siente el cuerpo, y vibrar daría a entender que hubo contacto.
  \item \textbf{Un único punto de salida (\cs{Haptica}):} toda la vibración del combate pasa por una clase estática que llama a \cs{OVRInput.SetControllerVibration} y avisa con un evento. Así se pudo verificar cada patrón registrando sus llamadas con marca de tiempo, y el registro de sesión cuenta las vibraciones.
\end{itemize}

Sobre RNF-02 (respuesta háptica en menos de 15 ms): la vibración se pide en la misma llamada en que se detecta el golpe, sin esperar ningún fotograma. La física corre a 50 Hz, así que la detección puede llegar hasta 20 ms después del contacto real; la optimización del paso de la física a la tasa de refresco del visor y la medición de la latencia forman parte del Entregable~5, según el cronograma.

\subsection{Registro de Métricas por Sesión (RF-21)}
\label{sec:registro}
El script \cs{RegistroSesion} se crea solo al iniciar la aplicación, sobrevive a los cambios de escena y registra cada combate escuchando eventos que ya existían, sin intervenir en la lógica del juego (tabla~\ref{tab:registro}).

\begin{table}[H]
\centering
\caption{Datos registrados por combate.}
\label{tab:registro}
\tablita{|>{\raggedright\arraybackslash}p{4.2cm}|L|}{\hd{Grupo}&\hd{Datos}}{
Combate & Nivel, rival, hora, si es un reintento, resultado (victoria o derrota por KO o por salud, o abandonado), round alcanzado y segundos de combate sin pausas\\ \hline
Golpes del jugador & Lanzados, conectados (por tipo), bloqueados y esquivados por el rival, precisión y efectividad\\ \hline
Defensa del jugador & Ataques recibidos, conectados, bloqueados y esquivados, y porcentaje de bloqueo\\ \hline
Caídas & Derribos al rival, caídas del jugador y cuentas superadas\\ \hline
Rendimiento y otros & Salud final de ambos, número de vibraciones, FPS medios y porcentaje de fotogramas por debajo de 72 FPS durante los rounds (RNF-01)\\ \hline
Por sesión & Duración, dispositivo, versión, combates, reintentos y victorias\\
}
\end{table}

Los datos se guardan en el dispositivo: un JSON por sesión, reescrito al terminar cada combate, al salir de una escena y al pausar o cerrar la aplicación, y un CSV (\cs{combates.csv}) con una fila por combate de todas las sesiones, listo para una hoja de cálculo. En el Quest se sacan con \cs{adb pull}.

\begin{itemize}
  \item \textbf{Precisión} = golpes conectados / lanzados. Un golpe se cuenta como lanzado cuando la mano se aleja de la cabeza a más de 1,5 m/s o cuando llega a tocar al rival. Aparte se da la \textbf{efectividad} (conectados / golpes que llegaron al rival), que no depende de ese umbral.
  \item \textbf{Reintento:} un combate de un nivel que ya se jugó en la misma sesión. \textbf{Abandonado:} se sale de la escena antes del final.
  \item \textbf{FPS solo durante los rounds,} porque RNF-01 habla de ``durante el combate''.
  \item \textbf{Formato local} en lugar de Unity Analytics: el archivo no necesita cuenta ni red y basta para las pruebas de la semana 14. Cubre los indicadores HEART de \emph{Engagement} (duración, combates reintentados) y \emph{Task Success} (precisión, golpes bloqueados, cuentas superadas) de la sección~3.6.
\end{itemize}

\subsection{Verificación y Validación}
\label{sec:verif}
\textbf{Pruebas en el Editor.} Se realizaron en Play en Unity 2022 LTS, con menús contextuales de prueba en cada componente y llamadas directas a los componentes. Sin visor la cabeza queda a la altura del piso, así que para probar los ataques del rival se colocaron la cabeza y las manos a mano. La tabla~\ref{tab:pruebas} resume lo comprobado.

\begin{table}[H]
\centering
\caption{Pruebas realizadas.}
\label{tab:pruebas}
\tablita{|>{\raggedright\arraybackslash}p{3.4cm}|L|}{\hd{Área}&\hd{Resultado}}{
Rondas & Combate completo en tiempo real: 90, 20, 90, 20 y 90 s con campana. El cronómetro se detiene al pausar y con el jugador caído. Victoria por salud (70\,\% contra 67\,\%), derrota (20\,\% contra 60\,\%) y empate a favor del jugador (50\,\% contra 50\,\%).\\ \hline
Progresión & Partida del nivel 1 al 4 hasta volver al menú: el nivel superado sube en cada victoria y se carga la escena siguiente. Perder por salud o por K.O. vuelve al menú sin cambiarlo. El nivel 4 aparece bloqueado hasta ganar el 3.\\ \hline
Cuenta del réferi & Del 1 al 10 en 10,00 s (rival) y 10,01 s (jugador). El round que llega a cero no se cierra a mitad de cuenta. La pausa detiene la cuenta y la voz, y se reanuda desde el mismo punto.\\ \hline
Combos de Soda y Sandman & Combo completo con tres impactos; un golpe, una caída, la pausa o el fin del round lo cortan sin impactos fantasma; con la guardia arriba, tres bloqueos; sacando la cabeza, tres esquives sin daño.\\ \hline
Esquive de Sandman & Se aparta 0,350 m y vuelve con desvío 0; no sale del ring; no esquiva atacando, caído ni incorporándose; el enfriamiento se respeta (nunca dos esquives a menos de 1,01 s).\\ \hline
Háptica & Con golpes reales de \cs{Puno} contra las zonas de golpe y registrando \cs{Haptica}: esquive (0,15 durante 46 ms), conectado (0,70 durante 65 ms), bloqueado, derribo (tres pulsos de 1,0) y un tic por segundo en la cuenta.\\ \hline
Guardia grabada & Grabación simulada con la cabeza girada 30° y inclinada 10°: 31 muestras recortadas; el JSON se importa como asset con las mismas muestras; los guantes ya no bajan de la pose animada.\\ \hline
Registro de sesión & Combates ganados, abandonados, reintentados y perdidos quedan en el JSON y en el CSV con los valores esperados (por ejemplo, 3 derribos por victoria).\\ \hline
Modelo de Sandman & Con \cs{ApoyoEnSuelo}, el vértice más bajo queda a 0,000 m de la lona en la guardia y en el combo (corrección de 0,23 a 0,37 m).\\ \hline
Compilación & Sin errores ni warnings propios en el Editor ni al compilar los scripts para Android.\\
}
\end{table}

\textbf{Validación con el visor.} Con el visor Meta Quest se validaron los aspectos que no se pueden comprobar en el Editor: la campana y las voces del réferi ubicadas en 3D, la sensación de cada patrón háptico, la grabación de la guardia con los dos sticks y su reproducción en los rivales, el \emph{button mashing} con los botones reales, el comportamiento de los combos y del esquive y su balance, la pausa del Guardián durante un round y la escritura del registro de sesión en el dispositivo. Las mediciones de FPS y de latencia háptica (RNF-01 y RNF-02) son parte del Entregable~5, con las pruebas HEART.

\subsection{Resumen de Cumplimiento (Semana 13)}
\begin{table}[H]
\centering
\caption{Cumplimiento del Entregable 4.}
\tablita{|>{\raggedright\arraybackslash}p{3.8cm}|L|c|}{\hd{Actividad planificada}&\hd{Implementación en Unity / scripts C\#}&\hd{Estado}}{
IA de los niveles 3 y 4 & \cs{AtaqueOponente} (combos), \cs{EsquiveOponente}, perfiles de Soda Popinski y Mr.\ Sandman, modelo de Mr.\ Sandman con \cs{ApoyoEnSuelo} & \ok\\ \hline
Réferi con cuenta regresiva y button mashing & \cs{Arbitro}, \cs{CuerpoArbitro}, \cs{RecuperacionJugador}, voz 3D sintetizada y modelo 3D del réferi & \ok\\ \hline
Rondas de 3 $\times$ 90 s & \cs{ControlRondas} y campana 3D & \ok\\ \hline
Hand Grab Recorder en los bloqueos & \cs{GrabadoraGuardia}, \cs{GrabacionGuardia} y \cs{DefensaOponente} & \ok\\ \hline
Háptica diferenciada & \cs{Puno}, \cs{SaludJugador}, \cs{FeedbackSalud} y \cs{Haptica} & \ok\\ \hline
Registro de métricas de sesión (adelantado) & \cs{RegistroSesion} & \ok\\ \hline
Resultado verificable & Cuatro rivales jugables y ciclo de KO funcionando & Cumplido\\
}
\end{table}

\subsection{Alcance del Entregable 5}
Quedan para la semana 14: la IA de Iron Mike Tyson y el nivel 5 (el modelo riggeado ya está en el proyecto), la escena de resultados con el cinturón de campeón y las estadísticas (RF-18 y RF-19, que leerán los datos que ya guarda el registro de sesión), la optimización a 72--90 FPS y la latencia háptica menor a 15 ms (RNF-01 y RNF-02), y las pruebas de usabilidad HEART (RNF-03 y RNF-05).

% ============================================================
\section{Implementación del Entregable 5: Versión Final de la Semana 14 (100\,\% de Avance Acumulado)}
\label{sec:e5}
\subsection{Resumen del Hito y Resultado Verificable}
Según el cronograma (sección~\ref{sec:crono}), el Entregable~5 comprende la IA del nivel 5 (Iron Mike Tyson), la escena de resultados con el cinturón y las estadísticas, la optimización a 72--90 FPS con háptica en menos de 15 ms, las pruebas HEART y la corrección de errores. Su resultado verificable es: \textbf{versión final con los 5 niveles y resultados de las pruebas de usabilidad}. La primera mitad se cumplió: el jugador recorre los cinco niveles en una partida continua, ve sus estadísticas al terminar cada combate y recibe el cinturón al vencer a Tyson. La segunda depende del visor y de los usuarios: \pendiente{resultados de las pruebas de usabilidad y mediciones en Meta Quest 2}.

Además se corrigieron los problemas que aparecieron al probar el juego en el visor: la pausa del Guardián que saltaba estando dentro del área segura, los personajes que se deslizaban sin dar pasos y el gesto rígido de la cuenta del réferi (sección~\ref{sec:correcciones}).

\begin{table}[H]
\centering
\caption{Bloques de trabajo del Entregable 5.}
\tablita{|c|L|>{\raggedright\arraybackslash}p{2.6cm}|c|}{\hd{Bloque}&\hd{Contenido}&\hd{Requisitos}&\hd{Sección}}{
H & Iron Mike Tyson y el nivel 5 & RF-07, RF-08, RF-09 & \ref{sec:tyson}\\ \hline
I & Escena de resultados y cinturón de campeón & RF-18, RF-19 & \ref{sec:resultados}\\ \hline
J & Rendimiento y latencia háptica & RNF-01, RNF-02, RNF-08 & \ref{sec:rendimiento}\\ \hline
K & Pruebas HEART con usuarios & RNF-03, RNF-05, RF-21 & \ref{sec:heart}\\ \hline
L & Correcciones tras la prueba en el visor & RF-12, RF-20, RNF-04 & \ref{sec:correcciones}\\ \hline
M & Informe final y comparación con los bocetos & --- & \ref{sec:bocetos}\\
}
\end{table}

\subsection{Iron Mike Tyson y el Nivel 5 (RF-07, RF-08, RF-09)}
\label{sec:tyson}
El equipo riggeó el modelo de Tyson en Blender con los nombres de hueso de Mixamo, así que Unity reconoció solo los 22 huesos del avatar humanoide. La malla original tenía 31\,092 triángulos, el doble que la de Mr.\ Sandman; un script de Blender sin interfaz (\cs{Tools/reducir_tyson.py}) la reduce a 15\,000 conservando los pesos y el UV. El mismo script corrige un defecto del rig: unos 130 vértices de lo alto de cada hombro dependían del antebrazo y, al subir la guardia, salía una ``cuchilla'' desde la axila. Como Mr.\ Sandman, Tyson usa las animaciones de Soda Popinski retargeteadas por el avatar humanoide, y el componente \cs{ApoyoEnSuelo} levanta la cadera lo que se hunden las suelas (de 6 a 14~cm con esos clips).

Su perfil es el más exigente del juego (tabla~\ref{tab:tyson}): combo de cuatro golpes (gancho izquierdo, gancho derecho, barrido y gancho izquierdo) cada 2 a 3~s, a 1,5 veces la velocidad de los clips, con más bloqueo y esquive que Mr.\ Sandman.

\begin{table}[H]
\centering
\caption{Perfil de Iron Mike Tyson frente a Mr.\ Sandman.}
\label{tab:tyson}
\tablita{|>{\raggedright\arraybackslash}p{5cm}|C|C|}{\hd{Parámetro}&\hd{Mr.\ Sandman}&\hd{Iron Mike Tyson}}{
Salud & 175 & 250\\ \hline
Bloqueo (base $\to$ máximo) & 25\,\% $\to$ 55\,\% & 30\,\% $\to$ 65\,\%\\ \hline
Esquive & 45\,\%, 1,0 s & 50\,\%, 0,9 s\\ \hline
Velocidad de ataque & $\times$1,4 & $\times$1,5\\ \hline
Patrón & gancho, barrido, gancho & gancho, gancho, barrido, gancho\\ \hline
Espera entre ataques & 2,8 a 4 s & 2 a 3 s\\ \hline
Daño por golpe (gancho / barrido) & 8 / 12 & 9 / 14\\ \hline
Castigo por guardia baja & +100\,\% & +120\,\%\\
}
\end{table}

Con una espera tan corta, Tyson lanzaba combos al aire mientras aún caminaba hacia el jugador. Desde este entregable ningún rival ataca si la cabeza del jugador está fuera de su alcance: sigue acercándose y ataca en cuanto llega.

Para comprobar que la dificultad sigue subiendo (RF-08) se repitieron las dos medidas del Entregable~4: el modelo de peleas simuladas y una pelea en Play con un jugador programado, igual en los niveles 4 y 5 (tabla~\ref{tab:dificultad5}). Tyson hace un 59\,\% más de daño por segundo, aguanta un 19\,\% más y derriba al jugador programado, que a Sandman nunca le cae; aun así, en las cuatro partidas ganó el jugador.

\begin{table}[H]
\centering
\caption{Dificultad de Iron Mike Tyson frente a Mr.\ Sandman (estimaciones).}
\label{tab:dificultad5}
\tablita{|>{\raggedright\arraybackslash}p{6.4cm}|C|C|}{\hd{Medida}&\hd{Mr.\ Sandman}&\hd{Iron Mike Tyson}}{
Daño por segundo si todo conecta, guardia arriba (modelo) & 6,65 & 10,60\\ \hline
Golpes para ganarle, 1 cada 0,5 s (modelo) & 52,8 & 80,9\\ \hline
Pelea simulada: segundos de round (dos partidas) & 98 y 99 & 118 y 117\\ \hline
Pelea simulada: daño recibido por el jugador & 91 y 86 & 150 y 138\\ \hline
Pelea simulada: caídas del jugador & 0 y 0 & 1 y 1\\
}
\end{table}

Son estimaciones con un jugador que nunca baja la guardia: \pendiente{confirmar en el visor, con jugadores reales, que el combate contra Tyson es difícil pero se puede ganar}.

\begin{figure}[H]
  \centering
  \begin{subfigure}[t]{0.48\linewidth}
    \includegraphics[width=\linewidth]{11_nivel5_mike_tyson.png}
    \caption{Iron Mike Tyson en guardia.}
  \end{subfigure}\hfill
  \begin{subfigure}[t]{0.48\linewidth}
    \includegraphics[width=\linewidth]{12_menu_niveles_5.png}
    \caption{Menú de niveles con el nivel 5.}
  \end{subfigure}
  \caption{El nivel 5 y su acceso desde el menú (se desbloquea al ganar el nivel 4, RF-09).}
\end{figure}

\subsection{Resultados y Cinturón de Campeón (RF-18, RF-19)}
\label{sec:resultados}
Al terminar cualquier combate, ganado o perdido, el juego muestra el cartel de siempre y carga \cs{Scene_Results}: la misma arena, con una hoja en un Canvas del mundo (no pegado a la vista) que se maneja con el puntero láser. Los datos los guarda el registro de sesión al cerrar el combate, así que la escena no necesita estado nuevo.

\begin{table}[H]
\centering
\caption{Contenido de la hoja de resultados.}
\tablita{|>{\raggedright\arraybackslash}p{3.6cm}|L|}{\hd{Elemento}&\hd{Qué muestra}}{
Título & ¡VICTORIA!, DERROTA o ¡CAMPEÓN!\\ \hline
Subtítulo & Nivel, rival y cómo terminó (por ejemplo, ``Ganaste por K.O. en el round 2'')\\ \hline
Precisión & Porcentaje y golpes conectados de los lanzados\\ \hline
Golpes bloqueados & Ataques del rival que el jugador bloqueó, de los recibidos\\ \hline
Tiempo total & Minutos y rounds del combate, sin las pausas\\ \hline
Detalle & Golpes que el rival bloqueó o esquivó, y derribos de cada uno\\ \hline
Botones & SIGUIENTE NIVEL (victoria), REINTENTAR (derrota) o REPETIR (tras el nivel 5), y MENÚ\\
}
\end{table}

\textbf{Cinturón de campeón.} No se usó el modelo de un repositorio externo, porque no se pudo verificar su licencia: el equipo generó uno propio con un script de Blender (\cs{Tools/crear_cinturon.py}), de 1\,540 triángulos, sin texturas y con tres materiales (cuero, oro y esmalte). Al ganar a Tyson, el cinturón sube desde el suelo creciendo, gira y su oro pulsa de brillo, y la hoja cambia al texto de cierre del arco narrativo: ``Derrotaste a Iron Mike Tyson, el campeón invicto. El cinturón de campeón regional es tuyo.'' Es el desenlace de la pirámide de Freytag (figura~\ref{fig:freytag}).

\begin{figure}[H]
  \centering
  \begin{subfigure}[t]{0.32\linewidth}
    \includegraphics[width=\linewidth]{13_resultados_victoria.png}
    \caption{Victoria.}
  \end{subfigure}\hfill
  \begin{subfigure}[t]{0.32\linewidth}
    \includegraphics[width=\linewidth]{14_resultados_derrota.png}
    \caption{Derrota.}
  \end{subfigure}\hfill
  \begin{subfigure}[t]{0.32\linewidth}
    \includegraphics[width=\linewidth]{15_resultados_campeon.png}
    \caption{Campeón, con el cinturón.}
  \end{subfigure}
  \caption{La hoja de resultados en sus tres variantes.}
\end{figure}

\subsection{Rendimiento y Latencia Háptica (RNF-01, RNF-02, RNF-08)}
\label{sec:rendimiento}
Las cifras de FPS y de latencia del Editor no valen para el Quest, así que este bloque deja preparado todo lo necesario para que la sesión con el visor consista en instalar, jugar y sacar los datos.

\textbf{Compilación para Android.} El juego completo compila para Meta Quest~2 (IL2CPP, ARM64, Vulkan, OpenXR con Multiview) sin errores. Hay tres APK generados desde el mismo código: el de desarrollo, el mismo sin MSAA (para comparar sin recompilar) y uno sin opciones de desarrollo. El 94\,\% del paquete son las texturas de 2\,048~px de la arena.

\textbf{Física a 72 Hz.} El paso de física pasó de 0,02~s (50 Hz) a exactamente 1/72~s, un paso por fotograma del Quest~2: el impacto se detecta como mucho un fotograma tarde, y no dos. Se revisó todo lo que depende del paso; solo el suavizado de la velocidad del puño necesitó ajustarse para seguir siendo el mismo filtro en segundos. Con golpes por física idénticos a 50 y a 72 Hz, los resultados y las clasificaciones fueron los mismos, y la medida de la rapidez del puño quedó más estable.

\textbf{Latencia medida en el juego.} La latencia háptica tiene tres tramos. El juego mide en cada golpe los dos que dependen de él y los guarda por combate (media, percentil 95 y máximo):

\begin{table}[H]
\centering
\caption{Tramos de la latencia háptica.}
\tablita{|>{\raggedright\arraybackslash}p{3cm}|L|>{\raggedright\arraybackslash}p{3.4cm}|}{\hd{Tramo}&\hd{Qué es}&\hd{Cómo se obtiene}}{
1. Detección & Del contacto real del guante al paso de física que lo detecta: de media medio fotograma (unos 7 ms) & Estimado en cada golpe por la penetración del guante en la zona\\ \hline
2. Código & Del aviso de colisión a la orden de vibración & Medido con un cronómetro de alta resolución\\ \hline
3. Runtime y mando & Del runtime de OpenXR al motor del mando por Bluetooth & No medible desde el código\\
}
\end{table}

RNF-02 se informará como la suma de los tramos 1 y 2 medidos en el visor: \pendiente{latencia háptica en Meta Quest 2 (media y p95)}.

\textbf{Auditoría de coste por escena.} Las siete escenas comparten la arena, que pesa casi todo el fotograma: unos 99\,000 triángulos y 570 piezas de material visibles, 25 luces en tiempo real con sombra y unas 175 texturas de 2\,048~px. Con esa auditoría se ordenaron doce recortes posibles por coste y beneficio. Ninguno está aplicado: se deciden con los FPS del visor. Los tres primeros (quitar las sombras de las luces salvo la direccional, bajar un nivel de mipmaps en Android y fijar la foveación) se aplican y se deshacen con un clic desde el menú \emph{Tools > Boxeo > Rendimiento}, que guarda el estado anterior de cada ajuste y lo restaura exactamente.

\textbf{Cambios tras la revisión del equipo.} Las texturas de los guantes pasaron de 2\,048~px sin comprimir a 1\,024~px en ASTC (de 107 a 5~MB de memoria), se activó el MSAA 4$\times$ en la calidad de Android (los bordes se veían dentados) y la aplicación se llama ``Simulador de Boxeo RV'' en la biblioteca del visor.

\begin{table}[H]
\centering
\caption{Rendimiento medido en Meta Quest 2 (se completa en la sesión con visor).}
\tablita{|>{\raggedright\arraybackslash}p{4.2cm}|L|c|}{\hd{Medida}&\hd{Resultado}&\hd{Cumple}}{
FPS medio por nivel (RNF-01) & \pendiente{tabla de \cs{Tools/tabla_rendimiento.py}} & \pendiente{}\\ \hline
Fotogramas bajo 72 FPS (RNF-01) & \pendiente{porcentaje del peor nivel} & \pendiente{}\\ \hline
Latencia háptica, tramos 1 + 2 (RNF-02) & \pendiente{p95 en ms} & \pendiente{}\\ \hline
Sesión continua (RNF-08) & \pendiente{duración, cierres y pérdidas de seguimiento} & \pendiente{}\\ \hline
MSAA 4$\times$ y recortes aplicados & \pendiente{decisión tras medir} & ---\\
}
\end{table}

\subsection{Pruebas HEART (RNF-03, RNF-05)}
\label{sec:heart}
El protocolo de las pruebas está escrito para que cualquiera del equipo pueda conducirlas: de 5 a 8 participantes novatos (3 usos de RV o menos), unos 25 minutos cada uno, con una persona que guía y vigila la seguridad y otra que anota y cronometra. Cada participante tiene un código (P01, P02...) y no se guarda su nombre en ningún archivo.

\begin{itemize}
  \item \textbf{Sin instrucciones de control (RNF-03).} Solo se explica cómo usar el puntero del menú. Desde la primera campana del nivel 1 no se dice cómo golpear, bloquear ni levantarse; si el participante pregunta, se responde ``haz lo que harías en un ring''.
  \item \textbf{Adopción medida por el juego.} El registro de sesión guarda, desde la primera campana y sin contar las pausas, los segundos hasta la primera guardia (las dos manos delante de la cara durante al menos 0,5 s, la misma condición que bloquea) y hasta el primer golpe que conecta. El cronómetro del anotador es el respaldo. Cumple quien hace las dos cosas en 120 s o menos.
  \item \textbf{Instrumentos.} Un cuestionario Likert de 1 a 5 (satisfacción inmersiva, ausencia de cinetosis y retención), una hoja de observación por participante y el registro de sesión del juego. Un script (\cs{Tools/tabla_heart.py}) cruza las tres fuentes por participante y calcula cada dimensión y los dos requisitos.
  \item \textbf{Retention adaptada.} No se pueden contar sesiones en la segunda y la tercera semana: se mide la intención de volver a jugar y si el participante sigue jugando cuando se le ofrece parar.
  \item \textbf{Observaciones dirigidas.} Si se entiende que hay que pulsar A, B, X o Y en la cuenta, si se nota el roce háptico cuando el rival esquiva y si los combos se pueden defender.
\end{itemize}

\begin{table}[H]
\centering
\caption{Resultados de las pruebas HEART (se completa tras las pruebas).}
\tablita{|>{\raggedright\arraybackslash}p{2.4cm}|>{\raggedright\arraybackslash}p{5cm}|L|}{\hd{Dimensión}&\hd{Indicador}&\hd{Resultado}}{
Happiness & Satisfacción inmersiva (media 1--5) & \pendiente{media}\\ \hline
Happiness & Ausencia de cinetosis, RNF-05 (mínimo 4) & \pendiente{media y si cumple}\\ \hline
Engagement & Duración de sesión, combates y reintentos & \pendiente{medias}\\ \hline
Adoption & Guardia y golpe en 2 min, RNF-03 (80\,\%) & \pendiente{porcentaje de novatos y si cumple}\\ \hline
Retention & Volvería a jugar; continuación voluntaria & \pendiente{porcentajes}\\ \hline
Task Success & Bloqueos, precisión, cuentas superadas & \pendiente{porcentajes}\\
}
\end{table}

\pendiente{lista priorizada de problemas encontrados en las pruebas y su corrección}.

\subsection{Correcciones tras la Prueba en el Visor (RF-12, RF-20, RNF-04)}
\label{sec:correcciones}
Al probar el juego en el Meta Quest se encontraron tres problemas, que se corrigieron:

\begin{itemize}
  \item \textbf{La pausa del Guardián saltaba estando dentro del área.} El área del Guardián que entrega el sistema está centrada en el centro del área dibujada, mientras que el juego mide la cabeza desde donde estaba el visor al arrancar. Si el jugador no empezaba justo en el centro, o recentraba la vista, el juego creía que estaba fuera y mostraba ``VUELVE AL ÁREA SEGURA''. Ahora el área se convierte al mismo sistema de coordenadas que la cabeza, se vuelve a leer al recentrar y al redibujar el Guardián, hace falta estar fuera un cuarto de segundo para pausar, y si el jugador pulsa REANUDAR estando ``fuera'' el juego deja de fiarse de esa geometría (el Guardián del propio Quest sigue avisando). Así un fallo nunca deja al jugador atrapado en la pausa.
  \item \textbf{Locomoción con joystick.} Un componente de depuración permitía moverse y girar con los sticks, lo que RNF-04 prohíbe. Ahora solo funciona en el Editor; en el visor el desplazamiento es únicamente físico.
  \item \textbf{Personajes que se deslizaban.} Ningún personaje tiene animación de caminar, así que al acercarse, esquivar o volver a su esquina se deslizaban sobre la lona. Un componente nuevo (\cs{PasosProcedurales}) clava cada pie en el suelo con cinemática inversa y, cuando el cuerpo se aleja o gira, da un paso en arco hasta su sitio, en cualquier dirección. Con el cuerpo quieto manda la animación.
  \item \textbf{Réferi.} Su gesto de cuenta era una bajada lineal del brazo recto. Ahora el brazo se mueve como un muelle amortiguado (baja rápido, se pasa un poco y se asienta), el antebrazo llega con retraso y el codo queda algo doblado, y el torso cabecea con cada número. Además se desplaza acelerando y frenando, y en la cuenta del rival \textbf{espera a que el rival esté tendido} para decir el primer número, en lugar de empezar mientras aún cae (unos 2 s después de la caída).
\end{itemize}

\pendiente{validar en el visor la pausa del Guardián, los pasos y el gesto del réferi}.

\subsection{Registro de Métricas (RF-21, Segunda Mitad)}
\label{sec:registro5}
El registro de sesión del Entregable~4 se amplió sin cambiar la lógica del juego: escucha eventos que ya existían. Por combate guarda ahora también la latencia háptica (tramos 1 y 2, con media, percentil 95 y máximo) y el paso de física con que se jugó, y para la adopción de RNF-03 el tiempo hasta la primera guardia y hasta el primer golpe conectado, por combate y por sesión. Las columnas nuevas van siempre al final del CSV: un archivo de una versión anterior solo amplía su cabecera, y sus filas viejas siguen siendo válidas.

\subsection{Comparación de los Bocetos con lo Implementado}
\label{sec:bocetos}
Esta comparación prepara el Entregable~6. Las tablas~\ref{tab:boceto-combate} y~\ref{tab:boceto-ko} recorren los elementos numerados de los bocetos de las figuras~\ref{fig:boceto-combate} y~\ref{fig:boceto-ko}, y la figura~\ref{fig:bocetos-real} muestra las mismas situaciones en el juego, desde los ojos del jugador.

\begin{table}[H]
\centering
\caption{Boceto de combate (figura~\ref{fig:boceto-combate}) frente a lo implementado.}
\label{tab:boceto-combate}
\tablita{|c|>{\raggedright\arraybackslash}p{2.6cm}|L|}{\hd{N.º}&\hd{Elemento}&\hd{Implementado}}{
1 & Esquive activo & Apartar la cabeza de la línea del golpe anula el daño (RF-06); sin vibración, porque ya lo siente el cuerpo\\ \hline
2 & Oponente IA & Cinco rivales con perfil propio, combos, bloqueo y esquive (secciones~\ref{sec:ia} y~\ref{sec:tyson})\\ \hline
3 & Tus guantes & Guantes 1:1 sobre los mandos, con texturas propias comprimidas para el Quest\\ \hline
4 & Salud & Viñeta que oscurece los bordes según la salud y parpadeo por debajo de 20 de salud (figura~\ref{fig:bocetos-real}a)\\ \hline
5 & Pulso cardíaco & Latido desde el 60\,\% de salud, de 60 a 90 pulsaciones por minuto y más fuerte cuanto menos salud (figura~\ref{fig:latido})\\ \hline
6 & Réferi & Réferi 3D con voz espacial que vigila la pelea a un lado y cuenta las caídas\\ \hline
7 & Vibración & Un patrón distinto por tipo de golpe, bloqueo, esquive, derribo y golpe recibido (sección~\ref{sec:haptica})\\
}
\end{table}

\begin{table}[H]
\centering
\caption{Boceto del estado de knockout (figura~\ref{fig:boceto-ko}) frente a lo implementado.}
\label{tab:boceto-ko}
\tablita{|c|>{\raggedright\arraybackslash}p{2.6cm}|L|}{\hd{N.º}&\hd{Elemento}&\hd{Implementado}}{
1 & Salud crítica & En la lona la viñeta llega a su máximo y la vista parpadea; con cada pulsación se aclara (figura~\ref{fig:bocetos-real}b)\\ \hline
2 & Pulso cardíaco & El latido acompaña la cuenta: empieza acelerado y se calma con cada pulsación\\ \hline
3 & Réferi & Se pone delante del jugador, a un lado del rival, cuenta del 1 al 10 en voz alta y marca cada número con el brazo, como en el boceto\\ \hline
4 & Button mashing & Pulsar A, B, X o Y levanta al jugador; cada pulsación vibra más fuerte y el cartel ``¡ARRIBA! PULSA A B X Y'' sobre el réferi lo indica, dentro del mundo y no como HUD\\
}
\end{table}

\begin{figure}[H]
  \centering
  \begin{subfigure}[t]{0.48\linewidth}
    \includegraphics[width=\linewidth]{16_salud_baja_combate.png}
    \caption{Combate con poca salud: los bordes se oscurecen.}
  \end{subfigure}\hfill
  \begin{subfigure}[t]{0.48\linewidth}
    \includegraphics[width=\linewidth]{17_cuenta_con_vineta.png}
    \caption{Knockout: viñeta al máximo, réferi contando y cartel de los botones.}
  \end{subfigure}
  \caption{Las situaciones de los bocetos en el juego (campo de visión de un ojo, recortado de la captura).}
  \label{fig:bocetos-real}
\end{figure}

\begin{figure}[H]
  \centering
  \begin{tikzpicture}[x=0.11cm,y=0.07cm,font=\small]
    \draw[->] (0,50) -- (105,50) node[below left]{salud (\%)};
    \draw[->] (0,50) -- (0,98) node[above right]{pulsaciones por minuto};
    \foreach \x in {0,20,40,60,80,100} \draw (\x,50) -- (\x,48.5) node[below]{\scriptsize\x};
    \foreach \y in {60,70,80,90} \draw (0,\y) -- (-1.5,\y) node[left]{\scriptsize\y};
    \draw[very thick,rojo,domain=0:60,samples=40] plot (\x,{60+30*((1-\x/60)^2*(3-2*(1-\x/60)))});
    \draw[dashed,gray] (60,50) -- (60,95);
    \node[align=left,gray] at (80,80) {\scriptsize sin latido\\[-2pt]\scriptsize por encima del 60\,\%};
  \end{tikzpicture}
  \caption{Ritmo del latido según la salud del jugador (el volumen sube igual, del 25\,\% al 100\,\%).}
  \label{fig:latido}
\end{figure}

\textbf{Diferencias con los bocetos.} El boceto del knockout mostraba solo el réferi con la mano en alto; en el juego el réferi también acompaña la pelea durante el round y cuenta al rival junto a él. El cartel ``¡ARRIBA! PULSA A B X Y'' no estaba en el boceto: se añadió porque, sin ningún texto, nada indicaba qué botón pulsar; está anclado al réferi dentro del mundo, así que respeta la regla de no usar HUD. El resto de los elementos se implementó como se diseñó.

\subsection{Verificación y Validación}
\label{sec:verif5}
\textbf{Pruebas en el Editor.} Se hicieron en Play con el Meta XR Simulator, con un arnés que fija la cabeza y las manos del jugador cada fotograma y, para los golpes, mueve el guante por trayectorias realistas contra las zonas de golpe, de modo que pasan por la física igual que en el visor.

\begin{table}[H]
\centering
\caption{Pruebas del Entregable 5 en el Editor.}
\label{tab:pruebas5}
\tablita{|>{\raggedright\arraybackslash}p{3.6cm}|L|}{\hd{Área}&\hd{Resultado}}{
Nivel 5 y progresión & Del menú al nivel 5 y de ahí a resultados; el nivel 5 bloqueado hasta ganar el 4\\ \hline
Resultados & Las tres variantes de la hoja y sus botones en los cinco niveles\\ \hline
Detección de impactos (RNF-07) & 70 golpes por física que atraviesan la zona de la cabeza, de 1 a 12 m/s y con desvíos al azar: 68 detectados (97\,\%); todos hasta 8 m/s y 9 de 10 a 10 y a 12 m/s. Los 21 roces del borde también se detectaron\\ \hline
Área segura (RNF-06) & Con un área de 1,5 $\times$ 1,5 m y la cabeza quieta en el centro, el rival se acerca y ataca con normalidad; a 0,6 m del centro sigue el combate, a 0,7 m (5 cm del borde) se pausa con el aviso y al volver se reanuda. Con un área equivocada, pulsar REANUDAR deja de usarla\\ \hline
Locomoción (RNF-04) & Ninguna escena tiene teletransporte; la locomoción con joystick se desactiva fuera del Editor\\ \hline
Adopción (RNF-03) & Guardia subida a los 3 s y golpe por física a los 5 s tras la campana: el registro guarda 3,03 y 5,22 s; tras un combate abandonado, el valor de sesión suma su tiempo\\ \hline
Cuenta del rival & El primer número llega 2,1 s después de la caída, con el rival ya tendido\\ \hline
Pasos & Rivales y réferi dan un paso cada 0,3 a 0,4 s mientras se desplazan; los pies no se hunden ni se elevan con el tiempo\\ \hline
Compilación & Los tres APK compilan sin errores; el paquete y los dispositivos declarados no cambian\\
}
\end{table}

\textbf{Validación con el visor.} \pendiente{FPS, latencia y sesión de 30 minutos según el protocolo de rendimiento; pausa del Guardián, pasos y gesto del réferi en el visor; pruebas HEART con usuarios}.

\subsection{Resumen de Cumplimiento (Semana 14)}
\begin{table}[H]
\centering
\caption{Cumplimiento del Entregable 5.}
\tablita{|>{\raggedright\arraybackslash}p{3.8cm}|L|>{\raggedright\arraybackslash}p{2.4cm}|}{\hd{Actividad planificada}&\hd{Implementación en Unity / scripts C\#}&\hd{Estado}}{
IA del nivel 5 & Perfil de Tyson, modelo reducido con \cs{ApoyoEnSuelo}, \cs{AtaqueOponente} sin ataques fuera de alcance & \ok\\ \hline
Resultados, cinturón y estadísticas & \cs{PantallaResultados}, \cs{CinturonCampeon} y \cs{Scene_Results} & \ok\\ \hline
Optimización (72--90 FPS, háptica < 15 ms) & Física a 1/72 s, latencia medida, auditoría y recortes con menú & \pendiente{medir en el visor}\\ \hline
Pruebas HEART & Protocolo, cuestionario, hoja de observación, \cs{Tools/tabla_heart.py} y adopción en el registro & \pendiente{pruebas con usuarios}\\ \hline
Corrección de errores & Pausa del Guardián, joystick, \cs{PasosProcedurales} y réferi & \ok\\ \hline
Resultado verificable & Versión final con los 5 niveles y resultados de las pruebas de usabilidad & \pendiente{resultados HEART}\\
}
\end{table}


\end{document}
